\chapter{Unit 5: Semiconductor Physics and Solid State Devices}

\section{Section 1: 100 Chapter-Wise Multiple-Choice Questions (MCQs)}

\subsection*{Part I: Foundational Level (Questions 1 to 50)}

\mcqitem{1}{The Wiedemann-Franz Law states that the ratio of thermal conductivity $K$ to electrical conductivity $\sigma$ for metals is:}{Directly proportional to absolute temperature $T$ ($K/\sigma = L T$)}{Inversely proportional to temperature $T$}{Independent of temperature}{Zero at all temperatures}
\mcqitem{2}{The classical Drude electrical conductivity $\sigma$ of a free electron gas with density $n$, relaxation time $\tau$, and electron mass $m$ is:}{$\sigma = \frac{n e^2 \tau}{m}$}{$\sigma = \frac{n e \tau}{m^2}$}{$\sigma = \frac{m}{n e^2 \tau}$}{$\sigma = n e \tau$}
\mcqitem{3}{At absolute zero temperature ($T = 0\text{ K}$), the Fermi-Dirac distribution function $f(E)$ for $E < E_F$ is:}{$1.0$ (completely occupied)}{$0.5$}{$0$}{$\infty$}
\mcqitem{4}{At absolute zero temperature ($T = 0\text{ K}$), the Fermi-Dirac distribution function $f(E)$ for $E > E_F$ is:}{$0$ (completely empty)}{$1.0$}{$0.5$}{$-1.0$}
\mcqitem{5}{At any finite temperature $T > 0\text{ K}$, the probability of occupancy of an energy state precisely at the Fermi level $E = E_F$ is:}{$0.5$ ($50\%$)}{$1.0$}{$0$}{$0.25$}
\mcqitem{6}{In an intrinsic (pure) semiconductor, the Fermi level $E_{Fi}$ lies:}{Essentially in the middle of the forbidden energy bandgap ($E_g/2$)}{Inside the conduction band}{Inside the valence band}{At the vacuum level}
\mcqitem{7}{The Kronig-Penney model explains the origin of energy bands in solids by analyzing electron motion in a:}{1D periodic square-well lattice potential representing crystalline atomic ion cores}{Completely uniform constant potential}{Coulomb potential of a single isolated atom}{Random chaotic potential}
\mcqitem{8}{In a direct bandgap semiconductor (such as GaAs or InP):}{The conduction band minimum and valence band maximum occur at the exact same crystal wavevector ($k = 0$)}{The band extrema occur at different $k$-values}{The bandgap is zero}{Photons cannot be emitted}
\mcqitem{9}{In an indirect bandgap semiconductor (such as Silicon or Germanium):}{Electron-hole recombination requires both a photon and a lattice phonon to conserve energy and crystal momentum simultaneously}{Photons can never be absorbed}{Effective mass is zero}{Energy bands do not exist}
\mcqitem{10}{Why is Gallium Arsenide (GaAs) used for LEDs and semiconductor lasers, whereas Silicon (Si) is not?}{GaAs is a direct bandgap semiconductor with high radiative efficiency, while Si is an indirect bandgap material where non-radiative phonon relaxation dominates}{Silicon is an insulator}{GaAs has zero electrical resistance}{Silicon cannot be doped}
\mcqitem{11}{The Law of Mass Action in non-degenerate semiconductors states that under thermal equilibrium:}{$n \cdot p = n_i^2$}{$n + p = n_i$}{$n / p = n_i$}{$n \cdot p = N_D N_A$}
\mcqitem{12}{An n-type semiconductor is created by doping pure Silicon (tetravalent, Group IV) with:}{Pentavalent donor impurity atoms from Group V (e.g., Phosphorus, Arsenic, Antimony)}{Trivalent acceptor impurity atoms from Group III (e.g., Boron, Gallium)}{Noble gas atoms}{Gold atoms only}
\mcqitem{13}{A p-type semiconductor is created by doping pure Silicon with:}{Trivalent acceptor impurity atoms from Group III (e.g., Boron, Aluminum, Gallium, Indium)}{Pentavalent donor atoms}{Iron atoms}{Oxygen molecules}
\mcqitem{14}{In an n-type semiconductor, the donor energy level $E_d$ lies:}{Just below the bottom of the conduction band $E_c$ ($\Delta E \approx 0.01\text{--}0.05\text{ eV}$)}{Just above the top of the valence band $E_v$}{In the exact middle of the bandgap}{Deep inside the valence band}
\mcqitem{15}{In a p-type semiconductor, the acceptor energy level $E_a$ lies:}{Just above the top of the valence band $E_v$ ($\Delta E \approx 0.01\text{--}0.05\text{ eV}$)}{Just below the conduction band}{At the center of the bandgap}{Inside the conduction band}
\mcqitem{16}{As the temperature of an extrinsic n-type semiconductor increases to very high levels ($T > 500\text{ K}$), the semiconductor behaves as:}{An intrinsic semiconductor (intrinsic carrier generation $n_i(T)$ overwhelms doping $N_D$)}{A superconductor}{A perfect dielectric}{A ferromagnetic metal}
\mcqitem{17}{The effective mass $m^*$ of an electron in a crystal band with energy dispersion relation $E(k)$ is given by:}{$m^* = \frac{\hbar^2}{\frac{d^2 E}{dk^2}}$}{$m^* = \hbar^2 \frac{d^2 E}{dk^2}$}{$m^* = \frac{1}{\hbar}\frac{dE}{dk}$}{$m^* = \frac{\hbar k}{E}$}
\mcqitem{18}{Drift velocity $v_d$ of charge carriers in an applied electric field $E$ is related to mobility $\mu$ by:}{$v_d = \mu E$}{$v_d = \mu / E$}{$v_d = \mu E^2$}{$v_d = E / \mu$}
\mcqitem{19}{The total electrical conductivity $\sigma$ of a semiconductor with electron mobility $\mu_n$ and hole mobility $\mu_p$ is:}{$\sigma = q (n \mu_n + p \mu_p)$}{$\sigma = q (n \mu_n - p \mu_p)$}{$\sigma = \frac{q}{n \mu_n + p \mu_p}$}{$\sigma = q n p (\mu_n + \mu_p)$}
\mcqitem{20}{Why is the electron mobility $\mu_n$ always greater than the hole mobility $\mu_p$ in Silicon ($\mu_n \approx 1350\text{ cm}^2/\text{V}\cdot\text{s}$ vs $\mu_p \approx 480\text{ cm}^2/\text{V}\cdot\text{s}$)?}{Electrons move in the conduction band with a smaller effective mass ($m_e^* < m_h^*$) compared to holes moving in the valence band}{Holes have negative charge}{Electrons are not affected by phonons}{Silicon contains no holes}
\mcqitem{21}{The Einstein relation connecting diffusion coefficient $D$ and carrier mobility $\mu$ at temperature $T$ is:}{$\frac{D_n}{\mu_n} = \frac{D_p}{\mu_p} = \frac{k_B T}{q} = V_T$}{$\frac{D}{\mu} = \frac{q}{k_B T}$}{$D \cdot \mu = \frac{k_B T}{q}$}{$\frac{D}{\mu} = k_B T q$}
\mcqitem{22}{The thermal voltage $V_T = \frac{k_B T}{q}$ at room temperature ($T = 300\text{ K}$) is approximately equal to:}{$25.9\text{ mV} \approx 26\text{ mV}$}{$1.0\text{ V}$}{$0.7\text{ V}$}{$120\text{ mV}$}
\mcqitem{23}{The Hall Effect occurs when a current-carrying semiconductor sample is placed in a transverse magnetic field, generating a transverse electric field due to the:}{Magnetic Lorentz force $\vec{F} = q(\vec{v}_d \times \vec{B})$ deflecting charge carriers to the lateral faces}{Coulomb repulsion between ions}{Thermal expansion of the semiconductor}{Photoelectric effect}
\mcqitem{24}{The Hall coefficient $R_H$ for an n-type semiconductor where electrons dominate is given by:}{$R_H = -\frac{1}{n q}$}{$R_H = +\frac{1}{n q}$}{$R_H = -n q$}{$R_H = 0$}
\mcqitem{25}{A positive Hall coefficient ($R_H > 0$) experimentally confirms that the semiconductor sample is:}{p-type (majority charge carriers are positive holes)}{n-type}{A perfect insulator}{A superconductor}
\mcqitem{26}{The Hall Voltage $V_H$ developed across a semiconductor slab of thickness $w$ carrying current $I$ in magnetic field $B$ is:}{$V_H = \frac{R_H I B}{w}$}{$V_H = R_H I B w$}{$V_H = \frac{I B}{R_H w}$}{$V_H = \frac{R_H w}{I B}$}
\mcqitem{27}{Hall mobility $\mu_H$ is related to Hall coefficient $R_H$ and electrical conductivity $\sigma$ by:}{$\mu_H = |R_H| \sigma$}{$\mu_H = \frac{|R_H|}{\sigma}$}{$\mu_H = \frac{\sigma}{|R_H|}$}{$\mu_H = |R_H| \sigma^2$}
\mcqitem{28}{In a p-n junction under thermal equilibrium (zero external bias), the net electric current across the junction is:}{Strictly zero ($I_{\text{drift}} + I_{\text{diffusion}} = 0$)}{Infinite}{Positive forward current}{Equal to saturation current $I_0$}
\mcqitem{29}{The built-in potential barrier $V_{bi}$ across a step-graded p-n junction at temperature $T$ is given by:}{$V_{bi} = \frac{k_B T}{q}\ln\left(\frac{N_A N_D}{n_i^2}\right)$}{$V_{bi} = \frac{k_B T}{q}\ln\left(\frac{n_i^2}{N_A N_D}\right)$}{$V_{bi} = \frac{q}{k_B T}\ln\left(\frac{N_A N_D}{n_i^2}\right)$}{$V_{bi} = \frac{k_B T}{q}\left(\frac{N_A + N_D}{n_i}\right)$}
\mcqitem{30}{The depletion region (space-charge layer) of a p-n junction consists of:}{Immobile ionized donor ions ($N_D^+$) on the n-side and immobile ionized acceptor ions ($N_A^-$) on the p-side, devoid of free mobile carriers}{High concentrations of free electrons and holes}{Pure neutral silicon atoms only}{Liquid electrolyte}
\mcqitem{31}{When a p-n junction is Forward Biased (positive terminal connected to p-side):}{The potential barrier height decreases ($V_{bi} - V_a$) and depletion layer width narrows, allowing large diffusion current}{The depletion layer widens}{The built-in potential increases}{Current drops to zero}
\mcqitem{32}{When a p-n junction is Reverse Biased (positive terminal connected to n-side):}{The potential barrier increases ($V_{bi} + V_R$) and depletion width widens, allowing only tiny minority carrier reverse saturation current $I_0$}{A massive forward current flows}{The barrier height drops to zero}{The bandgap vanishes}
\mcqitem{33}{The Shockley ideal diode equation relating current $I$ to applied voltage $V$ is:}{$I = I_0 \left(e^{qV/\eta k_BT} - 1\right)$}{$I = I_0 e^{qV/\eta k_BT}$}{$I = I_0 \left(e^{-qV/k_BT} + 1\right)$}{$I = \frac{V}{R_0}$}
\mcqitem{34}{The photovoltaic effect in a Solar Cell generates electricity through:}{Photogeneration of electron-hole pairs by incident photons ($h\nu \ge E_g$) and separation of carriers by the built-in electric field of a p-n junction}{Heating of the silicon wafer to boil water}{Chemical oxidation of metal electrodes}{Piezoelectric crystal deformation}
\mcqitem{35}{The Fill Factor ($FF$) of a solar cell with open-circuit voltage $V_{oc}$, short-circuit current $I_{sc}$, and maximum power point $(V_m, I_m)$ is:}{$FF = \frac{V_m I_m}{V_{oc} I_{sc}} = \frac{P_{\text{max}}}{V_{oc} I_{sc}}$}{$FF = \frac{V_{oc} I_{sc}}{V_m I_m}$}{$FF = \frac{V_m}{V_{oc}} + \frac{I_m}{I_{sc}}$}{$FF = V_{oc} I_{sc}$}
\mcqitem{36}{The power conversion efficiency $\eta$ of a solar cell illuminated by optical input power $P_{\text{in}}$ is:}{$\eta = \frac{P_{\text{max}}}{P_{\text{in}}} = \frac{V_{oc} I_{sc} FF}{P_{\text{in}}}$}{$\eta = \frac{P_{\text{in}}}{P_{\text{max}}}$}{$\eta = \frac{V_{oc}}{I_{sc} P_{\text{in}}}$}{$\eta = FF \times P_{\text{in}}$}
\mcqitem{37}{A Photodiode is operated in which electrical bias regime to achieve high speed and linearity?}{Reverse Bias}{Forward Bias}{Zero Bias (open circuit)}{High AC Voltage}
\mcqitem{38}{A Light Emitting Diode (LED) emits spontaneous light when operated in:}{Forward Bias}{Reverse Bias}{Breakdown regime}{Zero Bias}
\mcqitem{39}{The depletion width $W$ of a step-graded p-n junction with doping $N_A, N_D$ under applied voltage $V_a$ is proportional to:}{$W \propto \sqrt{V_{bi} - V_a}$}{$W \propto (V_{bi} - V_a)^2$}{$W \propto \frac{1}{V_{bi} - V_a}$}{$W \propto (V_{bi} - V_a)$}
\mcqitem{40}{In an asymmetric one-sided step junction ($p^+$-$n$ where $N_A \gg N_D$):}{The depletion region extends almost entirely into the lightly doped n-side ($W \approx x_n$)}{The depletion region extends equally into both sides}{The depletion region lies exclusively on the $p^+$ side}{No depletion region exists}
\mcqitem{41}{In intrinsic Silicon at $T = 300\text{ K}$, intrinsic carrier concentration is $n_i \approx 1.5 \times 10^{10}\text{ cm}^{-3}$. If Silicon is doped with $N_D = 10^{16}\text{ cm}^{-3}$ Arsenic atoms, the minority hole concentration $p$ is:}{$2.25 \times 10^4\text{ cm}^{-3}$}{$1.5 \times 10^6\text{ cm}^{-3}$}{$10^{16}\text{ cm}^{-3}$}{$2.25 \times 10^{10}\text{ cm}^{-3}$}
\mcqitem{42}{For the doped Silicon sample in Question 41, the shift of the Fermi level $(E_F - E_{Fi})$ above the intrinsic level is ($k_BT = 0.0259\text{ eV}$):}{$0.347\text{ eV}$}{$0.026\text{ eV}$}{$0.550\text{ eV}$}{$1.120\text{ eV}$}
\mcqitem{43}{A semiconductor sample of thickness $w = 0.5\text{ mm}$ carrying current $I = 10\text{ mA}$ in magnetic field $B = 0.5\text{ T}$ produces a Hall voltage $V_H = 5.0\text{ mV}$. The Hall coefficient $R_H$ is:}{$5.0 \times 10^{-4}\text{ m}^3/\text{C}$}{$5.0 \times 10^{-3}\text{ m}^3/\text{C}$}{$2.5 \times 10^{-2}\text{ m}^3/\text{C}$}{$1.0 \times 10^{-5}\text{ m}^3/\text{C}$}
\mcqitem{44}{For the sample in Question 43, the majority carrier concentration $n$ is:}{$1.25 \times 10^{22}\text{ m}^{-3} = 1.25 \times 10^{16}\text{ cm}^{-3}$}{$6.25 \times 10^{18}\text{ m}^{-3}$}{$1.25 \times 10^{14}\text{ cm}^{-3}$}{$3.00 \times 10^{15}\text{ cm}^{-3}$}
\mcqitem{45}{If the electrical conductivity of the sample in Question 43 is $\sigma = 270\text{ S/m}$, the Hall mobility $\mu_H$ is:}{$0.135\text{ m}^2/\text{V}\cdot\text{s} = 1350\text{ cm}^2/\text{V}\cdot\text{s}$}{$2700\text{ cm}^2/\text{V}\cdot\text{s}$}{$480\text{ cm}^2/\text{V}\cdot\text{s}$}{$500\text{ cm}^2/\text{V}\cdot\text{s}$}
\mcqitem{46}{In a Silicon p-n junction with $N_A = 10^{17}\text{ cm}^{-3}$ and $N_D = 10^{16}\text{ cm}^{-3}$ ($n_i = 1.5 \times 10^{10}\text{ cm}^{-3}$), the built-in potential $V_{bi}$ at $300\text{ K}$ is:}{$0.754\text{ V}$}{$0.550\text{ V}$}{$1.120\text{ V}$}{$0.350\text{ V}$}
\mcqitem{47}{A silicon solar cell has $V_{oc} = 0.60\text{ V}$, $I_{sc} = 3.0\text{ A}$, and $FF = 0.78$. The maximum electrical power output $P_{\text{max}}$ is:}{$1.404\text{ W}$}{$1.800\text{ W}$}{$2.340\text{ W}$}{$0.780\text{ W}$}
\mcqitem{48}{If the solar cell in Question 47 has an active area $A = 100\text{ cm}^2 = 0.01\text{ m}^2$ under standard $1\text{ Sun}$ illumination ($P_{\text{opt}} = 1000\text{ W/m}^2$), its efficiency $\eta$ is:}{$14.04\%$}{$18.00\%$}{$7.80\%$}{$23.40\%$}
\mcqitem{49}{The temperature dependence of carrier mobility due to acoustic phonon (lattice) scattering follows:}{$\mu_{\text{lattice}} \propto T^{-3/2}$}{$\mu_{\text{lattice}} \propto T^{+3/2}$}{$\mu_{\text{lattice}} \propto T^0$}{$\mu_{\text{lattice}} \propto T^{-1}$}
\mcqitem{50}{The temperature dependence of carrier mobility due to ionized impurity scattering follows:}{$\mu_{\text{impurity}} \propto T^{+3/2}$}{$\mu_{\text{impurity}} \propto T^{-3/2}$}{$\mu_{\text{impurity}} \propto 1/T$}{$\mu_{\text{impurity}} \propto e^{-T}$}

\subsection*{Part II: Intermediate Analytical Level (Questions 51 to 80)}

\mcqitem{51}{According to Matthiessen's Rule, the total carrier mobility $\mu$ combining lattice scattering ($\mu_L$) and impurity scattering ($\mu_I$) is:}{$\frac{1}{\mu} = \frac{1}{\mu_L} + \frac{1}{\mu_I}$}{$\mu = \mu_L + \mu_I$}{$\mu = \sqrt{\mu_L \mu_I}$}{$\mu = \mu_L - \mu_I$}
\mcqitem{52}{The intrinsic carrier concentration $n_i$ of Silicon as a function of temperature $T$ scales predominantly as:}{$n_i(T) \propto T^{3/2} \exp\left(-\frac{E_g}{2 k_B T}\right)$}{$n_i(T) \propto \exp\left(-\frac{E_g}{k_B T}\right)$}{$n_i(T) \propto T^3$}{$n_i(T) \propto \frac{1}{T}$}
\mcqitem{53}{If the energy bandgap of Silicon is $E_g = 1.12\text{ eV}$ and Germanium is $E_g = 0.67\text{ eV}$, at room temperature Germanium has:}{A much higher intrinsic carrier concentration $n_i$ and lower resistivity than Silicon}{Lower carrier concentration than Silicon}{Zero intrinsic carriers}{Identical resistivity to Silicon}
\mcqitem{54}{The position of the intrinsic Fermi level $E_{Fi}$ relative to the bandgap midpoint $\frac{E_c + E_v}{2}$ is given by:}{$E_{Fi} = \frac{E_c + E_v}{2} + \frac{3}{4}k_B T \ln\left(\frac{m_h^*}{m_e^*}\right)$}{$E_{Fi} = \frac{E_c + E_v}{2} - \frac{3}{4}k_B T \ln\left(\frac{m_h^*}{m_e^*}\right)$}{$E_{Fi} = \frac{E_c - E_v}{2}$}{$E_{Fi} = E_c - k_B T$}
\mcqitem{55}{If hole effective mass $m_h^*$ exceeds electron effective mass $m_e^*$ ($m_h^* > m_e^*$), as temperature $T$ increases the intrinsic Fermi level $E_{Fi}$:}{Shifts upward toward the conduction band}{Shifts downward toward the valence band}{Remains exactly static at absolute center}{Enters the core band}
\mcqitem{56}{In degenerate semiconductors (very heavily doped, $N_D > 10^{19}\text{ cm}^{-3}$):}{The Fermi level $E_F$ is driven inside the conduction band ($E_F > E_c$) and the material exhibits metallic conduction}{The bandgap becomes infinite}{Current stops flowing}{The material becomes an insulator}
\mcqitem{57}{The Zener breakdown in a p-n junction is caused by:}{Direct quantum mechanical tunneling of valence electrons on the p-side into the empty conduction band on the n-side under intense electric field ($> 10^6\text{ V/cm}$)}{Impact ionization by accelerated carriers}{Melting of the wire leads}{Thermal runaway}
\mcqitem{58}{The Avalanche breakdown in a p-n junction is caused by:}{Impact ionization where accelerated carriers gain enough kinetic energy to knock valence electrons into the conduction band, initiating an electron avalanche}{Direct quantum tunneling}{Laser excitation}{Chemical decomposition}
\mcqitem{59}{The temperature coefficient of breakdown voltage is:}{Negative for Zener breakdown and positive for Avalanche breakdown}{Positive for Zener and negative for Avalanche}{Positive for both}{Negative for both}
\mcqitem{60}{The junction capacitance $C_j$ of a reverse-biased step p-n junction varies with reverse voltage $V_R$ as:}{$C_j \propto \frac{1}{\sqrt{V_{bi} + V_R}}$}{$C_j \propto \sqrt{V_{bi} + V_R}$}{$C_j \propto (V_{bi} + V_R)^2$}{$C_j = \text{constant}$}
\mcqitem{61}{A Varactor Diode utilizes which electrical property of a p-n junction for RF voltage-controlled tuning?}{Voltage-dependent reverse junction capacitance $C_j(V_R)$}{Forward diffusion conductance}{Optical emission}{Zener resistance}
\mcqitem{62}{The diffusion capacitance $C_d$ of a p-n junction under Forward Bias is caused by:}{The storage of minority carrier excess charge distribution injected across the junction}{Immobile space charge ions}{The dielectric thickness of the casing}{Magnetic induction of the leads}
\mcqitem{63}{In Silicon, the electron diffusion length $L_n$ is related to electron diffusion coefficient $D_n$ and minority carrier lifetime $\tau_n$ by:}{$L_n = \sqrt{D_n \tau_n}$}{$L_n = D_n \tau_n$}{$L_n = \frac{D_n}{\tau_n}$}{$L_n = \sqrt{\frac{D_n}{\tau_n}}$}
\mcqitem{64}{Given $D_n = 35\text{ cm}^2/\text{s}$ and lifetime $\tau_n = 10\,\mu\text{s} = 10^{-5}\text{ s}$, the electron diffusion length $L_n$ is:}{$1.87 \times 10^{-2}\text{ cm} = 187\,\mu\text{m}$}{$350\,\mu\text{m}$}{$18.7\,\mu\text{m}$}{$1.87\text{ mm}$}
\mcqitem{65}{The Haynes-Shockley experiment is an important semiconductor characterization technique that directly measures:}{Minority carrier drift mobility $\mu$ and diffusion coefficient $D$}{The energy bandgap $E_g$}{The mass of the silicon atom}{The crystal lattice constant}
\mcqitem{66}{In the Hall effect, the Hall angle $\theta_H$ between total current density $\vec{J}$ and total electric field $\vec{E}$ is given by:}{$\tan \theta_H = \mu_H B$}{$\tan \theta_H = \frac{\mu_H}{B}$}{$\sin \theta_H = \mu_H B^2$}{$\cos \theta_H = \mu_H B$}
\mcqitem{67}{For an intrinsic semiconductor where electrons and holes both contribute to the Hall effect, the Hall coefficient $R_H$ is:}{$R_H = \frac{1}{q}\frac{p \mu_p^2 - n \mu_n^2}{(p \mu_p + n \mu_n)^2}$}{$R_H = \frac{1}{q(n+p)}$}{$R_H = 0$}{$R_H = \frac{\mu_n - \mu_p}{q}$}
\mcqitem{68}{The Quantum Hall Effect observed in 2D electron gases at cryogenic temperatures and high magnetic fields yields quantized Hall resistance:}{$R_H = \frac{h}{\nu e^2} \quad (\nu = 1, 2, 3,\dots)$}{$R_H = \frac{\nu h}{e}$}{$R_H = \frac{h^2}{\nu e}$}{$R_H = 0$}
\mcqitem{69}{The work function $\Phi_m$ of a metal and electron affinity $\chi_s$ of an n-type semiconductor determine the junction barrier. An ideal rectifying Schottky barrier forms when:}{$\Phi_m > \chi_s$ (for n-type semiconductor)}{$\Phi_m < \chi_s$}{$\Phi_m = \chi_s$}{$\chi_s = 0$}
\mcqitem{70}{An Ohmic contact (low-resistance non-rectifying contact) on an n-type semiconductor requires:}{$\Phi_m < \chi_s$ or extremely heavy surface doping ($n^+$ layer for field tunneling)}{$\Phi_m \gg \chi_s$}{An insulating oxide layer}{Reverse bias operation}
\mcqitem{71}{Schottky diodes have a major switching speed advantage over standard p-n junction diodes because:}{They are majority carrier devices with zero minority carrier storage time ($\tau_{rr} \approx 0$)}{They have infinite bandgap}{They operate without electrons}{They generate high reverse voltages}
\mcqitem{72}{The forward turn-on voltage of a Silicon Schottky diode is typically:}{$0.2\text{--}0.3\text{ V}$ (compared to $0.7\text{ V}$ for standard Si p-n diode)}{$0.7\text{ V}$}{$1.5\text{ V}$}{$5.0\text{ V}$}
\mcqitem{73}{In a Solar Cell under illumination, the short-circuit current $I_{sc}$ is directly proportional to:}{The incident optical illumination intensity $P_{\text{opt}}$}{The square of the illumination intensity}{The temperature only}{The load resistance}
\mcqitem{74}{The open-circuit voltage $V_{oc}$ of an ideal solar cell is given by:}{$V_{oc} = \frac{k_B T}{q}\ln\left(\frac{I_{sc}}{I_0} + 1\right)$}{$V_{oc} = \frac{I_{sc} R_s}{I_0}$}{$V_{oc} = \frac{k_B T}{q}\left(\frac{I_{sc}}{I_0}\right)$}{$V_{oc} = V_{bi}$}
\mcqitem{75}{As the operating temperature of a Silicon solar cell increases, its efficiency $\eta$:}{Decreases (due to exponential increase in reverse saturation current $I_0$ reducing $V_{oc}$)}{Increases linearly}{Remains perfectly constant}{Doubles every $10^\circ\text{C}$}
\mcqitem{76}{Anti-reflective coatings on solar cells (such as $\text{Si}_3\text{N}_4$ or $\text{TiO}_2$) with refractive index $n_{\text{ARC}} = \sqrt{n_{\text{Si}} n_{\text{air}}} \approx 1.9$ reduce bare silicon reflection from $\approx 35\%$ down to:}{$< 2\text{--}3\%$}{$25\%$}{$50\%$}{$0\%$ across all wavelengths}
\mcqitem{77}{A Solar Cell has $V_{oc} = 0.62\text{ V}$, $I_{sc} = 4.0\text{ A}$, and maximum power $P_{\text{max}} = 1.984\text{ W}$. The Fill Factor $FF$ is:}{$0.80$}{$0.70$}{$0.60$}{$0.90$}
\mcqitem{78}{The Shockley-Queisser theoretical efficiency limit for a single-junction Silicon solar cell under standard 1-Sun unconcentrated sunlight is approximately:}{$\approx 33.7\%$}{$100\%$}{$50.0\%$}{$15.0\%$}
\mcqitem{79}{In a Tunnel Diode (Esaki Diode), heavy degenerate doping on both p and n sides ($p^{++}$-$n^{++}$) creates a unique $I$-$V$ characteristic featuring:}{Negative Differential Resistance (NDR) in the forward bias region}{Infinite resistance}{Zero current in all regimes}{Strictly linear Ohm's law behavior}
\mcqitem{80}{The Negative Differential Resistance (NDR) of a tunnel diode is widely used in high-frequency engineering for:}{Microwave oscillators, high-speed pulse generators, and amplifiers}{Solar power conversion}{Battery charging}{Audio loudspeakers}

\subsection*{Part III: Advanced \& Numerical Problems (Questions 81 to 100)}

\mcqitem{81}{In the Kronig-Penney model, the dispersion relation is given by $P\frac{\sin(\alpha a)}{\alpha a} + \cos(\alpha a) = \cos(ka)$. The barrier strength parameter $P$ is defined as:}{$P = \frac{m V_0 b a}{\hbar^2}$}{$P = \frac{V_0}{E}$}{$P = \frac{\hbar^2 k^2}{2m}$}{$P = \frac{a}{b}$}
\mcqitem{82}{In the limit $P \to \infty$ in the Kronig-Penney dispersion relation, the energy spectrum reduces to:}{Discrete, atomic-like bound states $\alpha a = n\pi \implies E_n = \frac{n^2 h^2}{8ma^2}$}{A single continuous band}{Zero energy everywhere}{Negative infinity}
\mcqitem{83}{The first Brillouin zone for a 1D crystal lattice with lattice constant $a$ extends over the wavevector range:}{$-\frac{\pi}{a} \le k \le +\frac{\pi}{a}$}{$0 \le k \le \frac{\pi}{a}$}{$-\frac{2\pi}{a} \le k \le +\frac{2\pi}{a}$}{$-\infty < k < +\infty$}
\mcqitem{84}{At the edge of the first Brillouin zone ($k = \pm \pi/a$), the electron group velocity $v_g = \frac{1}{\hbar}\frac{dE}{dk}$ is:}{$0$ (standing wave formation via Bragg reflection)}{$c$}{$v_F$}{$\infty$}
\mcqitem{85}{The effective density of states in the conduction band $N_c$ for a 3D parabolic band is given by:}{$N_c = 2 \left(\frac{2\pi m_e^* k_B T}{h^2}\right)^{3/2}$}{$N_c = \left(\frac{m_e^* k_B T}{2\pi \hbar^2}\right)^{1/2}$}{$N_c = 2 \frac{k_B T}{h}$}{$N_c = \frac{n_i^2}{N_v}$}
\mcqitem{86}{The effective density of states in the valence band $N_v$ is similarly:}{$N_v = 2 \left(\frac{2\pi m_h^* k_B T}{h^2}\right)^{3/2}$}{$N_v = N_c \frac{m_e^*}{m_h^*}$}{$N_v = \sqrt{N_c}$}{$N_v = 0$}
\mcqitem{87}{In the freeze-out temperature regime of an extrinsic n-type semiconductor ($T < 100\text{ K}$):}{Thermal energy $k_BT$ is insufficient to ionize donor atoms, causing electrons to remain trapped in donor levels $E_d$}{All electrons enter the conduction band}{Resistivity drops to zero}{The semiconductor melts}
\mcqitem{88}{In the extrinsic saturation (exhaustion) regime ($150\text{ K} \le T \le 450\text{ K}$):}{All donor atoms are $100\%$ ionized, maintaining a constant electron concentration $n \approx N_D$ independent of temperature}{Carrier concentration drops exponentially}{Intrinsic carriers exceed donors}{Mobility becomes infinite}
\mcqitem{89}{The Early Effect (Base-Width Modulation) in a Bipolar Junction Transistor (BJT) is caused by:}{The widening of the reverse-biased collector-base depletion region narrowing the effective neutral base width $W_B$}{Base melting}{Thermal noise in emitter}{Zener tunneling in the base}
\mcqitem{90}{The Gunn Effect in III-V semiconductors (such as GaAs) arises from:}{Field-induced transfer of conduction electrons from a high-mobility lower central valley ($\Gamma$) to low-mobility higher satellite valleys ($L$)}{Acoustic phonon absorption}{Photoelectric ionization}{Nuclear magnetic resonance}
\mcqitem{91}{A Gunn Diode is widely used as a solid-state source of:}{Microwave and Millimeter-wave radiation ($10\text{ GHz to }100\text{ GHz}$)}{Visible red light}{High-voltage DC power}{Audio sound}
\mcqitem{92}{The IMPATT (Impact Ionization Avalanche Transit Time) diode generates high microwave power by combining:}{Avalanche multiplication and carrier transit-time delay to create a $180^\circ$ phase lag between current and voltage}{Direct bandgap emission}{Photo-conduction}{Superconducting flux pinning}
\mcqitem{93}{The flat-band voltage $V_{\text{FB}}$ of a Metal-Oxide-Semiconductor (MOS) structure is given by:}{$V_{\text{FB}} = \Phi_{\text{ms}} - \frac{Q_{\text{ox}}}{C_{\text{ox}}}$}{$V_{\text{FB}} = \Phi_{\text{ms}} + \frac{Q_{\text{ox}}}{C_{\text{ox}}}$}{$V_{\text{FB}} = \frac{Q_{\text{ox}}}{C_{\text{ox}}}$}{$V_{\text{FB}} = 0$}
\mcqitem{94}{In a MOSFET, strong inversion occurs when the surface potential $\psi_s$ reaches:}{$\psi_s = 2 \phi_F = 2 \left(\frac{k_BT}{q}\right)\ln\left(\frac{N_A}{n_i}\right)$}{$\psi_s = \phi_F$}{$\psi_s = 0$}{$\psi_s = V_{DD}$}
\mcqitem{95}{The subthreshold swing $S$ of an ideal MOSFET at room temperature ($300\text{ K}$) has a fundamental thermodynamic lower limit of:}{$S \ge \ln(10)\frac{k_BT}{q} \approx 60\text{ mV/decade}$}{$S = 0\text{ mV/decade}$}{$S = 10\text{ mV/decade}$}{$S = 100\text{ V/decade}$}
\mcqitem{96}{Quantum confinement in a 2D electron gas (2DEG) formed at an $\text{AlGaAs/GaAs}$ heterojunction interface leads to:}{Ultra-high electron mobility ($> 10^6\text{ cm}^2/\text{V}\cdot\text{s}$ at $4\text{ K}$) due to spatial separation of carriers from ionized donor dopants (Modulation Doping)}{Zero conductivity}{Insulating behavior}{Loss of Fermi level}
\mcqitem{97}{The Quantum Point Contact (QPC) demonstrates ballistic 1D transport where conductance is quantized in universal steps of:}{$G = n \left(\frac{2 e^2}{h}\right) \quad (n = 1, 2, 3,\dots)$}{$G = n \frac{e}{h}$}{$G = \frac{h}{2e^2}$}{$G = 0$}
\mcqitem{98}{In thermoelectric energy conversion, the figure of merit $ZT$ of a semiconductor is defined as:}{$ZT = \frac{S^2 \sigma}{\kappa} T$}{$ZT = \frac{S \sigma}{\kappa T}$}{$ZT = \frac{\kappa}{S^2 \sigma T}$}{$ZT = S \sigma \kappa T$}
\mcqitem{99}{Why are heavy-element semiconductors like Bismuth Telluride ($\text{Bi}_2\text{Te}_3$) and Silicon-Germanium ($\text{SiGe}$) preferred for thermoelectric generators?}{They exhibit high power factors ($S^2\sigma$) and extremely low lattice thermal conductivity ($\kappa_{\text{lattice}}$) due to heavy atomic mass phonon scattering}{They have infinite bandgap}{They are transparent}{They are magnetic}
\mcqitem{100}{The Hall coefficient for a high-purity p-type Silicon sample is measured as $R_H = +3.5 \times 10^{-2}\text{ m}^3/\text{C}$. The hole concentration $p$ is:}{$1.78 \times 10^{20}\text{ m}^{-3} = 1.78 \times 10^{14}\text{ cm}^{-3}$}{$3.50 \times 10^{16}\text{ cm}^{-3}$}{$1.00 \times 10^{18}\text{ cm}^{-3}$}{$5.00 \times 10^{12}\text{ cm}^{-3}$}

\newpage
\section{Section 2: Complete Answer Key with Detailed Rationales}

\subsection*{Master Answer Key (Questions 1 to 25)}

{\small
\noindent\begin{tabularx}{\textwidth}{|c|c|X|}
\hline
\textbf{Q\#} & \textbf{Ans} & \textbf{Technical Rationale \& Calculation} \\
\hline
\textbf{1} & \textbf{(a)} & The Wiedemann-Franz Law establishes that $K/\sigma = L T$, where Lorenz number $L = \frac{\pi^2 k_B^2}{3e^2} \approx 2.44 \times 10^{-8}\text{ W}\cdot\Omega/\text{K}^2$. \\ \hline
\textbf{2} & \textbf{(a)} & Drude free electron theory gives $\sigma = n e^2\tau / m$. \\ \hline
\textbf{3} & \textbf{(a)} & At $T = 0\text{ K}$, all energy states below the Fermi level $E_F$ are $100\%$ occupied ($f(E) = 1$). \\ \hline
\textbf{4} & \textbf{(a)} & At $T = 0\text{ K}$, all energy states above $E_F$ are completely unoccupied ($f(E) = 0$). \\ \hline
\textbf{5} & \textbf{(a)} & $f(E_F) = \frac{1}{1 + e^{(E_F - E_F)/k_BT}} = \frac{1}{1 + e^0} = \frac{1}{1 + 1} = 0.5$. \\ \hline
\textbf{6} & \textbf{(a)} & For equal effective masses, $E_{Fi} = \frac{E_c + E_v}{2}$, located precisely at the bandgap midpoint. \\ \hline
\textbf{7} & \textbf{(a)} & The Kronig-Penney model applies Bloch's theorem to a periodic array of potential barriers, generating allowed energy bands and forbidden gaps. \\ \hline
\textbf{8} & \textbf{(a)} & Direct band alignment permits rapid radiative electron-hole recombination with direct photon emission without phonon participation. \\ \hline
\textbf{9} & \textbf{(a)} & Because conduction minimum and valence maximum are displaced in $k$-space, radiative transitions require a phonon to satisfy $\Delta k \ne 0$. \\ \hline
\textbf{10} & \textbf{(a)} & Direct bandgap GaAs produces efficient light emission, whereas indirect bandgap Silicon releases recombination energy as thermal lattice heat. \\ \hline
\textbf{11} & \textbf{(a)} & The product of electron concentration $n$ and hole concentration $p$ is a constant dependent only on temperature: $n p = n_i^2(T)$. \\ \hline
\textbf{12} & \textbf{(a)} & Group V donor atoms donate extra conduction electrons, creating majority carrier electron concentration $n \approx N_D$. \\ \hline
\textbf{13} & \textbf{(a)} & Group III acceptor atoms accept electrons from the valence band, creating majority carrier hole concentration $p \approx N_A$. \\ \hline
\textbf{14} & \textbf{(a)} & Donor levels lie only tens of $\text{meV}$ below $E_c$, allowing room-temperature thermal energy ($k_BT \approx 26\text{ meV}$) to ionize nearly all donors. \\ \hline
\textbf{15} & \textbf{(a)} & Acceptor levels lie slightly above $E_v$, so valence electrons easily jump into $E_a$, leaving behind free mobile holes in the valence band. \\ \hline
\textbf{16} & \textbf{(a)} & At high temperatures, thermal band-to-band excitation causes $n_i(T) \gg N_D$, causing the material to transition into the intrinsic regime. \\ \hline
\textbf{17} & \textbf{(a)} & Effective mass reflects lattice force interactions: $m^* = \hbar^2 / (d^2E/dk^2)$. \\ \hline
\textbf{18} & \textbf{(a)} & Drift velocity is directly proportional to applied electric field: $v_d = \mu E$. \\ \hline
\textbf{19} & \textbf{(a)} & Total conductivity is the sum of electron and hole drift conductivities: $\sigma = q(n\mu_n + p\mu_p)$. \\ \hline
\textbf{20} & \textbf{(a)} & Conduction band curvature is sharper than valence band curvature, giving electrons smaller effective mass and higher mobility. \\ \hline
\textbf{21} & \textbf{(a)} & The Einstein relation establishes $\frac{D}{\mu} = \frac{k_BT}{q} = V_T \approx 25.9\text{ mV}$ at room temperature ($300\text{ K}$). \\ \hline
\textbf{22} & \textbf{(a)} & $V_T = \frac{(1.38 \times 10^{-23}\text{ J/K})(300\text{ K})}{1.602 \times 10^{-19}\text{ C}} \approx 0.02585\text{ V} \approx 25.9\text{ mV}$. \\ \hline
\textbf{23} & \textbf{(a)} & Moving carriers experience transverse Lorentz force $q(\vec{v}_d \times \vec{B})$, accumulating on side faces to produce the Hall voltage. \\ \hline
\textbf{24} & \textbf{(a)} & For n-type material, majority carriers are electrons ($q = -e$), giving negative Hall coefficient $R_H = -1/(ne)$. \\ \hline
\textbf{25} & \textbf{(a)} & A positive sign for $R_H = +1/(pq)$ proves that holes are the majority mobile charge carriers. \\ \hline
\end{tabularx}
}

\newpage

\subsection*{Master Answer Key (Questions 26 to 50)}

{\small
\noindent\begin{tabularx}{\textwidth}{|c|c|X|}
\hline
\textbf{Q\#} & \textbf{Ans} & \textbf{Technical Rationale \& Calculation} \\
\hline
\textbf{26} & \textbf{(a)} & Balancing electric and magnetic forces gives $V_H = E_H d = \frac{R_H J B}{\dots} \implies V_H = \frac{R_H I B}{w}$. \\ \hline
\textbf{27} & \textbf{(a)} & Combining $R_H = 1/(nq)$ and $\sigma = nq\mu$ yields the Hall mobility formula $\mu_H = |R_H|\sigma$. \\ \hline
\textbf{28} & \textbf{(a)} & At equilibrium, hole diffusion current is exactly balanced by hole drift current, and electron diffusion balances electron drift, giving zero net current. \\ \hline
\textbf{29} & \textbf{(a)} & Integrating Poisson's equation across the depletion region yields $V_{bi} = V_T \ln(N_A N_D / n_i^2)$. \\ \hline
\textbf{30} & \textbf{(a)} & Carrier diffusion leaves behind exposed uncompensated fixed donor ($N_D^+$) and acceptor ($N_A^-$) ionic space charge. \\ \hline
\textbf{31} & \textbf{(a)} & Forward bias opposes the built-in field, reducing the barrier height and allowing majority carriers to diffuse readily across the junction. \\ \hline
\textbf{32} & \textbf{(a)} & Reverse bias reinforces the built-in electric field, widening the space-charge width $W \propto \sqrt{V_{bi} + V_R}$. \\ \hline
\textbf{33} & \textbf{(a)} & The ideal diode current-voltage relationship is $I = I_0 (e^{qV/\eta k_BT} - 1)$, where $\eta$ is the ideality factor. \\ \hline
\textbf{34} & \textbf{(a)} & Photons create electron-hole pairs in the depletion region; the built-in field sweeps electrons to the n-side and holes to the p-side, producing photovoltage. \\ \hline
\textbf{35} & \textbf{(a)} & Fill Factor measures squareness of the $I$-$V$ curve: $FF = P_{\text{max}} / (V_{oc} I_{sc}) < 1.0$ (typically $0.7\text{--}0.85$). \\ \hline
\textbf{36} & \textbf{(a)} & Solar cell efficiency is the ratio of maximum electrical power output to total incident optical irradiance power: $\eta = (V_{oc} I_{sc} FF)/P_{\text{in}}$. \\ \hline
\textbf{37} & \textbf{(a)} & Reverse bias widens the depletion region, minimizes junction capacitance ($C_j \propto 1/W$), and accelerates carrier transit time. \\ \hline
\textbf{38} & \textbf{(a)} & Forward bias injects majority carriers across the junction, resulting in radiative electron-hole recombination. \\ \hline
\textbf{39} & \textbf{(a)} & From Poisson's equation, depletion width scales as $W = \sqrt{\frac{2\epsilon_s}{q}\left(\frac{1}{N_A} + \frac{1}{N_D}\right)(V_{bi} - V_a)}$. \\ \hline
\textbf{40} & \textbf{(a)} & Charge neutrality requires $q N_A x_p = q N_D x_n$. Since $N_A \gg N_D$, $x_p \ll x_n$, so depletion penetrates into the lightly doped region. \\ \hline
\textbf{41} & \textbf{(a)} & By Mass Action: $p = \frac{n_i^2}{n} \approx \frac{(1.5 \times 10^{10})^2}{10^{16}} = \frac{2.25 \times 10^{20}}{10^{16}} = 2.25 \times 10^4\text{ cm}^{-3}$. \\ \hline
\textbf{42} & \textbf{(a)} & $E_F - E_{Fi} = k_BT \ln\left(\frac{N_D}{n_i}\right) = 0.0259 \ln\left(\frac{10^{16}}{1.5 \times 10^{10}}\right) = 0.0259 \ln(6.67 \times 10^5) = 0.0259(13.41) \approx 0.347\text{ eV}$. \\ \hline
\textbf{43} & \textbf{(a)} & $R_H = \frac{V_H w}{I B} = \frac{(5 \times 10^{-3}\text{ V})(0.5 \times 10^{-3}\text{ m})}{(10 \times 10^{-3}\text{ A})(0.5\text{ T})} = \frac{2.5 \times 10^{-6}}{5 \times 10^{-3}} = 5.0 \times 10^{-4}\text{ m}^3/\text{C}$. \\ \hline
\textbf{44} & \textbf{(a)} & $n = \frac{1}{|R_H| q} = \frac{1}{(5.0 \times 10^{-4})(1.602 \times 10^{-19})} \approx 1.248 \times 10^{22}\text{ m}^{-3} = 1.25 \times 10^{16}\text{ cm}^{-3}$. \\ \hline
\textbf{45} & \textbf{(a)} & $\mu_H = |R_H|\sigma = (5.0 \times 10^{-4}\text{ m}^3/\text{C})(270\text{ S/m}) = 0.135\text{ m}^2/\text{V}\cdot\text{s} = 1350\text{ cm}^2/\text{V}\cdot\text{s}$. \\ \hline
\textbf{46} & \textbf{(a)} & $V_{bi} = 0.0259 \ln\left(\frac{10^{17} \times 10^{16}}{(1.5 \times 10^{10})^2}\right) = 0.0259 \ln\left(\frac{10^{33}}{2.25 \times 10^{20}}\right) = 0.0259 \ln(4.44 \times 10^{12}) = 0.0259(29.12) \approx 0.754\text{ V}$. \\ \hline
\textbf{47} & \textbf{(a)} & $P_{\text{max}} = V_{oc} \cdot I_{sc} \cdot FF = (0.60\text{ V})(3.0\text{ A})(0.78) = 1.404\text{ W}$. \\ \hline
\textbf{48} & \textbf{(a)} & $P_{\text{in}} = 1000\text{ W/m}^2 \times 0.01\text{ m}^2 = 10\text{ W}$. $\eta = \frac{P_{\text{max}}}{P_{\text{in}}} = \frac{1.404\text{ W}}{10\text{ W}} = 0.1404 = 14.04\%$. \\ \hline
\textbf{49} & \textbf{(a)} & As temperature increases, thermal lattice vibrations increase, causing electron-phonon collision rates to rise as $\mu \propto T^{-3/2}$. \\ \hline
\textbf{50} & \textbf{(a)} & Faster-moving thermal carriers spend less time near ionized dopant ions, reducing Coulomb deflection and increasing mobility as $\mu \propto T^{3/2}$. \\ \hline
\end{tabularx}
}

\newpage

\subsection*{Master Answer Key (Questions 51 to 75)}

{\small
\noindent\begin{tabularx}{\textwidth}{|c|c|X|}
\hline
\textbf{Q\#} & \textbf{Ans} & \textbf{Technical Rationale \& Calculation} \\
\hline
\textbf{51} & \textbf{(a)} & Scattering rates sum linearly: $\frac{1}{\tau_{\text{total}}} = \frac{1}{\tau_L} + \frac{1}{\tau_I} \implies \frac{1}{\mu} = \frac{1}{\mu_L} + \frac{1}{\mu_I}$. \\ \hline
\textbf{52} & \textbf{(a)} & $n_i = \sqrt{N_c N_v}e^{-E_g/2k_BT}$, and since $N_c, N_v \propto T^{3/2}$, $n_i \propto T^{3/2} e^{-E_g/(2k_BT)}$. \\ \hline
\textbf{53} & \textbf{(a)} & Because $E_g(\text{Ge}) < E_g(\text{Si})$, $e^{-E_g/2k_BT}$ is orders of magnitude larger for Ge ($n_i \approx 2.4 \times 10^{13}\text{ cm}^{-3}$ vs $1.5 \times 10^{10}\text{ cm}^{-3}$ in Si). \\ \hline
\textbf{54} & \textbf{(a)} & Setting $n = p \implies N_c e^{-(E_c-E_F)/k_BT} = N_v e^{-(E_F-E_v)/k_BT}$ gives $E_{Fi} = \frac{E_c+E_v}{2} + \frac{3}{4}k_BT\ln(m_h^*/m_e^*)$. \\ \hline
\textbf{55} & \textbf{(a)} & Because $\ln(m_h^*/m_e^*) > 0$, the positive term $\frac{3}{4}k_BT\ln(m_h^*/m_e^*)$ increases linearly with $T$, shifting $E_{Fi}$ upward. \\ \hline
\textbf{56} & \textbf{(a)} & Extreme donor doping forces $E_F$ above the conduction band edge $E_c$, causing the semiconductor to behave degenerately like a metal. \\ \hline
\textbf{57} & \textbf{(a)} & Zener breakdown occurs in heavily doped thin junctions ($V_Z < 5\text{ V}$) via direct quantum band-to-band tunneling. \\ \hline
\textbf{58} & \textbf{(a)} & In lightly doped wide junctions ($V_{\text{BR}} > 6\text{ V}$), carriers trigger avalanche multiplication through high-energy impact ionization. \\ \hline
\textbf{59} & \textbf{(a)} & Zener tunneling increases with temperature (bandgap narrows $\implies$ negative tempco); avalanche breakdown decreases with temperature (more phonon scattering $\implies$ positive tempco). \\ \hline
\textbf{60} & \textbf{(a)} & Because $C_j = \frac{\epsilon_s A}{W}$ and depletion width $W \propto \sqrt{V_{bi} + V_R}$, capacitance scales as $C_j \propto 1/\sqrt{V_{bi} + V_R}$. \\ \hline
\textbf{61} & \textbf{(a)} & Varactors operate as voltage-variable capacitors in reverse bias for RF frequency synthesizers and voltage-controlled oscillators (VCOs). \\ \hline
\textbf{62} & \textbf{(a)} & Forward bias injects minority carriers whose stored charge $Q = I\tau$ varies with voltage, creating diffusion capacitance $C_d = \frac{q I \tau}{k_BT}$. \\ \hline
\textbf{63} & \textbf{(a)} & Solving the steady-state continuity equation $\frac{d^2(\Delta n)}{dx^2} = \frac{\Delta n}{D_n\tau_n}$ defines diffusion length $L_n = \sqrt{D_n\tau_n}$. \\ \hline
\textbf{64} & \textbf{(a)} & $L_n = \sqrt{D_n \tau_n} = \sqrt{35 \times 10^{-5}} = \sqrt{3.5 \times 10^{-4}} \approx 0.0187\text{ cm} = 187\,\mu\text{m}$. \\ \hline
\textbf{65} & \textbf{(a)} & The Haynes-Shockley experiment injects a minority pulse and tracks its drift time and dispersion to measure $\mu$ and $D$ directly. \\ \hline
\textbf{66} & \textbf{(a)} & $\tan\theta_H = E_H / E_x = (R_H J B)/ (J/\sigma) = R_H \sigma B = \mu_H B$. \\ \hline
\textbf{67} & \textbf{(a)} & Two-carrier transport gives $R_H = \frac{1}{q}\frac{p\mu_p^2 - n\mu_n^2}{(p\mu_p + n\mu_n)^2}$. Since $\mu_n > \mu_p$, intrinsic Si has a slightly negative $R_H$. \\ \hline
\textbf{68} & \textbf{(a)} & Klaus von Klitzing discovered that 2D Hall resistance is quantized in integer steps of the von Klitzing constant $R_K = h/e^2 \approx 25812.807\,\Omega$. \\ \hline
\textbf{69} & \textbf{(a)} & For $\Phi_m > \chi_s$, electrons flow from semiconductor to metal until a built-in barrier $q\phi_{Bn} = \Phi_m - \chi_s$ forms, creating a rectifying Schottky diode. \\ \hline
\textbf{70} & \textbf{(a)} & Ohmic contacts have negligible barrier or are so heavily doped that electrons easily tunnel through the thin barrier. \\ \hline
\textbf{71} & \textbf{(a)} & Conduction is purely via majority electrons, eliminating minority carrier recombination delays and allowing switching in picoseconds. \\ \hline
\textbf{72} & \textbf{(a)} & Lower barrier heights in Schottky junctions yield lower forward voltage drops ($\approx 0.2\text{--}0.3\text{ V}$), reducing conduction power losses. \\ \hline
\textbf{73} & \textbf{(a)} & Photogeneration rate $G = \eta_{\text{opt}}\Phi_{\text{photon}}$ is linearly proportional to incident photon flux: $I_{sc} \approx I_L \propto P_{\text{opt}}$. \\ \hline
\textbf{74} & \textbf{(a)} & Setting net current $I = 0 = I_{sc} - I_0(e^{qV_{oc}/k_BT}-1)$ yields $V_{oc} = V_T \ln(I_{sc}/I_0 + 1)$. \\ \hline
\textbf{75} & \textbf{(a)} & Higher temperature dramatically increases intrinsic carrier density $n_i^2$, raising $I_0$ and reducing $V_{oc}$ by $\approx -2\text{ mV/}^\circ\text{C}$. \\ \hline
\end{tabularx}
}

\newpage

\subsection*{Master Answer Key (Questions 76 to 100)}

{\small
\noindent\begin{tabularx}{\textwidth}{|c|c|X|}
\hline
\textbf{Q\#} & \textbf{Ans} & \textbf{Technical Rationale \& Calculation} \\
\hline
\textbf{76} & \textbf{(a)} & Quarter-wavelength destructive interference coatings ($d = \lambda_0 / (4n_{\text{ARC}})$) eliminate front surface Fresnel reflections. \\ \hline
\textbf{77} & \textbf{(a)} & $FF = \frac{P_{\text{max}}}{V_{oc} I_{sc}} = \frac{1.984}{(0.62)(4.0)} = \frac{1.984}{2.48} = 0.80$. \\ \hline
\textbf{78} & \textbf{(a)} & Thermodynamic bandgap balance between unabsorbed sub-bandgap photons and thermalization losses limits single-junction cells to $\approx 33.7\%$. \\ \hline
\textbf{79} & \textbf{(a)} & Degenerate alignment of occupied conduction states with empty valence states allows tunneling, which drops as forward bias increases (NDR region). \\ \hline
\textbf{80} & \textbf{(a)} & NDR offsets circuit resistance losses, enabling microwave frequency oscillations in the gigahertz regime. \\ \hline
\textbf{81} & \textbf{(a)} & The dimensionless barrier strength $P = \frac{m V_0 b a}{\hbar^2}$ measures electron confinement; $P \to 0$ gives free electrons, $P \to \infty$ gives isolated atoms. \\ \hline
\textbf{82} & \textbf{(a)} & As $P \to \infty$, allowed energy bands shrink into infinitely sharp discrete atomic energy eigenvalues. \\ \hline
\textbf{83} & \textbf{(a)} & The fundamental Wigner-Seitz cell in reciprocal space (first Brillouin zone) spans from $-\pi/a$ to $+\pi/a$. \\ \hline
\textbf{84} & \textbf{(a)} & At band edges, incident and Bragg-reflected waves form standing waves, meaning net forward group velocity vanishes ($v_g = 0$). \\ \hline
\textbf{85} & \textbf{(a)} & Integrating density of states with Maxwell-Boltzmann tail gives $N_c = 2(2\pi m_e^* k_BT / h^2)^{3/2}$. \\ \hline
\textbf{86} & \textbf{(a)} & Similarly, the valence band effective density of states is $N_v = 2(2\pi m_h^* k_BT / h^2)^{3/2}$. \\ \hline
\textbf{87} & \textbf{(a)} & At cryogenic temperatures, $k_BT \ll (E_c - E_d)$, so electrons freeze back into localized donor impurity states. \\ \hline
\textbf{88} & \textbf{(a)} & Over normal operating temperatures, all dopants are fully ionized while intrinsic generation is negligible, keeping $n \approx N_D$ constant. \\ \hline
\textbf{89} & \textbf{(a)} & Increasing reverse collector voltage $V_{CB}$ widens the space-charge layer into the base, reducing neutral base width $W_B$. \\ \hline
\textbf{90} & \textbf{(a)} & The Ridley-Watkins-Hilsum (RWH) transferred-electron mechanism creates negative differential mobility, producing microwave Gunn oscillations. \\ \hline
\textbf{91} & \textbf{(a)} & Transit of high-field dipole domains generates continuous microwave power in radar systems and wireless transceivers. \\ \hline
\textbf{92} & \textbf{(a)} & Impact ionization ($90^\circ$ phase lag) plus drift transit delay ($90^\circ$) yields a total $180^\circ$ phase shift, generating dynamic negative resistance. \\ \hline
\textbf{93} & \textbf{(a)} & Flat-band voltage accounts for metal-semiconductor work function difference $\Phi_{\text{ms}}$ and fixed oxide charge: $V_{\text{FB}} = \Phi_{\text{ms}} - Q_{\text{ox}}/C_{\text{ox}}$. \\ \hline
\textbf{94} & \textbf{(a)} & Strong inversion begins when the surface electron concentration equals the bulk majority hole concentration: $\psi_s = 2\phi_F$. \\ \hline
\textbf{95} & \textbf{(a)} & Boltzmann thermal tail limits subthreshold gate swing to $S = \ln(10)V_T \approx 2.3026 \times 25.85\text{ mV} \approx 60\text{ mV/decade}$. \\ \hline
\textbf{96} & \textbf{(a)} & High Electron Mobility Transistors (HEMT) use modulation doping to separate electrons in GaAs from ionized donors in AlGaAs, eliminating impurity scattering. \\ \hline
\textbf{97} & \textbf{(a)} & Landauer formula for ballistic transport gives quantized conductance steps of $G_0 = \frac{2e^2}{h} \approx 77.48\,\mu\text{S} = \frac{1}{12.9\text{ k}\Omega}$. \\ \hline
\textbf{98} & \textbf{(a)} & Thermoelectric efficiency depends on Seebeck coefficient $S$, electrical conductivity $\sigma$, and thermal conductivity $\kappa$: $ZT = \frac{S^2 \sigma}{\kappa}T$. \\ \hline
\textbf{99} & \textbf{(a)} & Heavy atomic masses and complex unit cells scatter heat-carrying phonons, suppressing thermal conductivity $\kappa$ while preserving electrical conductivity $\sigma$. \\ \hline
\textbf{100} & \textbf{(a)} & $p = \frac{1}{R_H q} = \frac{1}{(0.035\text{ m}^3/\text{C})(1.602 \times 10^{-19}\text{ C})} \approx 1.783 \times 10^{20}\text{ m}^{-3} = 1.78 \times 10^{14}\text{ cm}^{-3}$. \\ \hline
\end{tabularx}
}


\newpage
\section{Section 3: 20 Comprehensive Theory Questions with Analytical Derivations}

\begin{theoryq}{1}
Explain the Classical Drude Free Electron Theory of Metals. Derive the expression for electrical conductivity $\sigma = \frac{n e^2 \tau}{m}$. State and derive the Wiedemann-Franz Law, and discuss the physical meaning of the Lorenz Number.
\end{theoryq}
\begin{theorysol}
\textbf{1. Assumptions of Drude Model}:
\begin{enumerate}
    \item A metal consists of positive lattice ion cores and a gas of free valence electrons of number density $n$.
    \item In the absence of an electric field, electrons move randomly with thermal velocity $\vec{v}_{\text{th}}$ ($\langle \vec{v}_{\text{th}} \rangle = 0$).
    \item Under an applied electric field $\vec{E}$, electrons experience an electrostatic force $\vec{F} = -e\vec{E}$ and accelerate between collisions.
    \item Collisions with lattice ions are instantaneous, resetting drift velocity to zero after average relaxation time $\tau$.
\end{enumerate}

\textbf{2. Derivation of Electrical Conductivity $\sigma$}:
The equation of motion of an electron under electric field $\vec{E}$ is:
\begin{equation}
m \frac{d\vec{v}_d}{dt} + \frac{m}{\tau}\vec{v}_d = -e\vec{E}
\end{equation}
In steady state ($\frac{d\vec{v}_d}{dt} = 0$):
\begin{equation}
\vec{v}_d = -\frac{e\tau}{m}\vec{E} = -\mu_n \vec{E}
\end{equation}
The resulting drift current density $\vec{J}$ is:
\begin{equation}
\vec{J} = -n e \vec{v}_d = -n e \left( -\frac{e\tau}{m}\vec{E} \right) = \left( \frac{n e^2 \tau}{m} \right)\vec{E}
\end{equation}
Comparing with Ohm's law $\vec{J} = \sigma \vec{E}$:
\begin{equation}
\mathbf{\sigma = \frac{n e^2 \tau}{m}}
\end{equation}

\textbf{3. Thermal Conductivity $K$ \& Wiedemann-Franz Law}:
From classical kinetic theory of gases, thermal conductivity carried by electron gas is:
\begin{equation}
K = \frac{1}{3} C_v v_{\text{th}} \lambda_{\text{mfp}} = \frac{1}{3}\left(\frac{3}{2}n k_B\right) v_{\text{th}} (v_{\text{th}}\tau) = \frac{1}{2} n k_B \tau v_{\text{th}}^2
\end{equation}
Using equipartition theorem $\frac{1}{2}m v_{\text{th}}^2 = \frac{3}{2}k_BT \implies v_{\text{th}}^2 = \frac{3k_BT}{m}$:
\begin{equation}
K = \frac{3}{2}\frac{n k_B^2 T \tau}{m}
\end{equation}
Taking the ratio of thermal to electrical conductivity:
\begin{equation}
\mathbf{\frac{K}{\sigma} = \frac{\frac{3}{2}\frac{n k_B^2 T \tau}{m}}{\frac{n e^2 \tau}{m}} = \frac{3}{2}\left(\frac{k_B}{e}\right)^2 T = L_{\text{classical}} T}
\end{equation}
In quantum Sommerfeld theory (using Fermi-Dirac statistics):
\begin{equation}
\mathbf{\frac{K}{\sigma} = \frac{\pi^2}{3}\left(\frac{k_B}{e}\right)^2 T = L T}
\end{equation}
where $L = \frac{\pi^2 k_B^2}{3e^2} \approx \mathbf{2.44 \times 10^{-8}\text{ W}\cdot\Omega/\text{K}^2}$ is the universal \textbf{Lorenz Number}.
\end{theorysol}
\vspace{4mm}

\begin{theoryq}{2}
Explain the physical origin of energy bands in solids using the Kronig-Penney Model. State Bloch's theorem and derive how periodic boundary conditions lead to allowed energy bands and forbidden gaps.
\end{theoryq}
\begin{theorysol}
\textbf{1. Bloch's Theorem}:
For an electron moving in a 1D periodic crystalline lattice potential $V(x + a) = V(x)$ of lattice constant $a$, the Schrödinger wavefunctions are modulated plane waves:
\begin{equation}
\mathbf{\psi_k(x) = e^{i k x} u_k(x)}
\end{equation}
where $u_k(x + a) = u_k(x)$ is a periodic function with the periodicity of the lattice, and $k$ is the crystal wavevector.

\textbf{2. Kronig-Penney Potential Model}:
The periodic crystal potential is modeled as an array of rectangular barrier wells:
\begin{equation}
V(x) = \begin{cases} 0 & 0 < x < a \quad (\text{well}) \\ V_0 & -b < x < 0 \quad (\text{barrier}) \end{cases}
\end{equation}
Period is $d = a + b$.

\textbf{3. Dispersion Relation}:
Applying boundary conditions on $\psi$ and $\frac{d\psi}{dx}$ at boundaries, and taking the Dirac delta-function limit ($V_0 \to \infty, b \to 0$ such that $V_0 b = \text{finite}$), yields the fundamental \textbf{Kronig-Penney Relation}:
\begin{equation}
\mathbf{P \frac{\sin(\alpha a)}{\alpha a} + \cos(\alpha a) = \cos(k a)}
\end{equation}
where $\alpha = \frac{\sqrt{2mE}}{\hbar}$ and $P = \frac{m V_0 b a}{\hbar^2}$ is the dimensionless barrier strength.

\textbf{4. Formation of Allowed Bands \& Forbidden Gaps}:
Because the right side $\cos(ka)$ is bounded strictly between $-1$ and $+1$:
\begin{equation}
\mathbf{-1 \le P \frac{\sin(\alpha a)}{\alpha a} + \cos(\alpha a) \le +1}
\end{equation}
\begin{itemize}
    \item \textbf{Allowed Energy Bands}: Values of energy $E = \frac{\hbar^2\alpha^2}{2m}$ where the function lies inside $[-1, +1]$. In these ranges, real solutions for $k$ exist.
    \item \textbf{Forbidden Energy Gaps ($E_g$)}: Values of energy where the magnitude of the function exceeds $1.0$ ($|f(\alpha a)| > 1$). No real wavevector $k$ exists, creating forbidden energy gaps where electron transmission is impossible.
\end{itemize}
\end{theorysol}
\vspace{4mm}

\begin{theoryq}{3}
Derive the expressions for electron concentration $n$ and hole concentration $p$ in an intrinsic semiconductor. Obtain the Law of Mass Action $n \cdot p = n_i^2$ and derive the position of the intrinsic Fermi level $E_{Fi}$.
\end{theoryq}
\begin{theorysol}
\textbf{1. Electron Concentration in Conduction Band $n$}:
The electron density is obtained by integrating the density of states $g_c(E)$ multiplied by the Fermi-Dirac probability $f(E)$ from $E_c$ to $\infty$:
\begin{equation}
n = \int_{E_c}^\infty g_c(E) f(E) dE = \frac{1}{2\pi^2}\left(\frac{2m_e^*}{\hbar^2}\right)^{3/2} \int_{E_c}^\infty \frac{\sqrt{E - E_c}}{1 + e^{(E - E_F)/k_BT}} dE
\end{equation}
In non-degenerate semiconductors ($E_c - E_F \ge 3k_BT$), using the Maxwell-Boltzmann tail $f(E) \approx e^{-(E - E_F)/k_BT}$:
\begin{equation}
\mathbf{n = N_c e^{-(E_c - E_F)/k_BT}}
\end{equation}
where $N_c = 2\left(\frac{2\pi m_e^* k_BT}{h^2}\right)^{3/2}$ is the effective density of states in the conduction band.

\textbf{2. Hole Concentration in Valence Band $p$}:
Integrating vacant states $[1 - f(E)]$ in the valence band:
\begin{equation}
\mathbf{p = N_v e^{-(E_F - E_v)/k_BT}}
\end{equation}
where $N_v = 2\left(\frac{2\pi m_h^* k_BT}{h^2}\right)^{3/2}$ is the effective density of states in the valence band.

\textbf{3. Law of Mass Action}:
Multiplying $n$ and $p$:
\begin{align*}
n \cdot p &= \left(N_c e^{-(E_c - E_F)/k_BT}\right)\left(N_v e^{-(E_F - E_v)/k_BT}\right) = N_c N_v e^{-(E_c - E_v)/k_BT} \\
\mathbf{n \cdot p} &= \mathbf{N_c N_v e^{-E_g / k_BT} = n_i^2}
\end{align*}
where intrinsic carrier concentration is $\mathbf{n_i = \sqrt{N_c N_v} e^{-E_g / (2k_BT)}}$.

\textbf{4. Intrinsic Fermi Level $E_{Fi}$}:
In an intrinsic semiconductor, $n = p = n_i$:
\begin{align*}
N_c e^{-(E_c - E_{Fi})/k_BT} &= N_v e^{-(E_{Fi} - E_v)/k_BT} \\
e^{2E_{Fi}/k_BT} &= \frac{N_v}{N_c} e^{(E_c + E_v)/k_BT} = \left(\frac{m_h^*}{m_e^*}\right)^{3/2} e^{(E_c + E_v)/k_BT}
\end{align*}
Taking natural logarithms:
\begin{equation}
\frac{2E_{Fi}}{k_BT} = \frac{E_c + E_v}{k_BT} + \frac{3}{2}\ln\left(\frac{m_h^*}{m_e^*}\right) \implies \mathbf{E_{Fi} = \frac{E_c + E_v}{2} + \frac{3}{4}k_BT \ln\left(\frac{m_h^*}{m_e^*}\right)}
\end{equation}
\end{theorysol}
\vspace{4mm}

\begin{theoryq}{4}
Derive the complete mathematical expressions for the Hall Voltage $V_H$, Hall Electric Field $E_H$, Hall Coefficient $R_H$, and Hall Mobility $\mu_H$. Explain four vital practical engineering applications of the Hall Effect.
\end{theoryq}
\begin{theorysol}
\textbf{1. Physical Setup \& Force Balance}:
Consider a rectangular semiconductor slab of width $w$ (along $y$), thickness $d$ (along $z$), and length $L$ (along $x$).
A direct current $I_x$ is passed along the $+x$ direction, and a uniform magnetic field $B_z$ is applied along the $+z$ direction.
\begin{enumerate}
    \item Charge carriers moving with drift velocity $\vec{v}_d = v_x \hat{i}$ experience magnetic Lorentz force:
    \begin{equation}
    \vec{F}_m = q (\vec{v}_d \times \vec{B}) = q (v_x \hat{i} \times B_z \hat{k}) = -q v_x B_z \hat{j}
    \end{equation}
    \item Deflected carriers accumulate on side faces, setting up a transverse electrostatic Hall electric field $\vec{E}_H = E_y \hat{j}$.
    \item In steady state, electrostatic force balances magnetic Lorentz force ($F_e + F_m = 0$):
    \begin{equation}
    q E_H = q v_x B_z \implies \mathbf{E_H = v_x B_z}
    \end{equation}
\end{enumerate}

\textbf{2. Derivation of Hall Coefficient $R_H$ \& Hall Voltage $V_H$}:
Current density is $J_x = n q v_x \implies v_x = \frac{J_x}{n q}$.
Substituting $v_x$ into the electric field balance:
\begin{equation}
E_H = \left(\frac{J_x}{n q}\right) B_z = \left(\frac{1}{n q}\right) J_x B_z
\end{equation}
Defining the \textbf{Hall Coefficient $R_H$}:
\begin{equation}
\mathbf{R_H = \frac{E_H}{J_x B_z} = \frac{1}{n q}} \implies \begin{cases} R_H = -\frac{1}{n e} & (\text{n-type}) \\ R_H = +\frac{1}{p e} & (\text{p-type}) \end{cases}
\end{equation}
The induced transverse Hall Voltage $V_H$ across slab width $w$ is:
\begin{equation}
\mathbf{V_H = E_H \cdot w = R_H J_x B_z w = R_H \left(\frac{I_x}{w \cdot d}\right) B_z w = \frac{R_H I_x B_z}{d}}
\end{equation}
where $d$ is slab thickness along the magnetic field.

\textbf{3. Hall Mobility $\mu_H$}:
Using conductivity $\sigma = n q \mu$:
\begin{equation}
\mathbf{\mu_H = |R_H| \sigma = \left(\frac{1}{n q}\right)(n q \mu) = \mu}
\end{equation}

\textbf{4. Four Key Engineering Applications}:
\begin{enumerate}
    \item \textbf{Determination of Semiconductor Type}: Sign of $V_H$ ($R_H > 0 \implies \text{p-type}$; $R_H < 0 \implies \text{n-type}$).
    \item \textbf{Measurement of Carrier Density}: $n = 1/(|R_H|e)$.
    \item \textbf{Measurement of Carrier Mobility}: $\mu = |R_H|\sigma$.
    \item \textbf{Magnetic Field Sensors \& Brushless DC Motor Commutators}: Linear Hall probes accurately measure magnetic flux densities from $\mu\text{T}$ to $> 10\text{ T}$.
\end{enumerate}

\end{theorysol}
\vspace{4mm}

\begin{theoryq}{5}
Derive the step-by-step expression for the built-in potential barrier $V_{bi}$ and depletion layer width $W$ across an abrupt p-n junction from Poisson's equation.
\end{theoryq}
\begin{theorysol}
\textbf{1. Poisson's Equation Formulation}:
Consider an abrupt p-n junction with acceptor doping $N_A$ on p-side ($-x_p \le x \le 0$) and donor doping $N_D$ on n-side ($0 \le x \le x_n$).
\begin{equation}
\frac{d^2 V}{dx^2} = -\frac{\rho(x)}{\epsilon_s} = \begin{cases} +\frac{q N_A}{\epsilon_s} & -x_p \le x \le 0 \\ -\frac{q N_D}{\epsilon_s} & 0 \le x \le x_n \end{cases}
\end{equation}

\textbf{2. Integrating for Electric Field $E(x) = -\frac{dV}{dx}$}:
Applying boundary conditions $E(-x_p) = 0$ and $E(x_n) = 0$:
\begin{equation}
E(x) = \begin{cases} -\frac{q N_A}{\epsilon_s}(x + x_p) & -x_p \le x \le 0 \\ -\frac{q N_D}{\epsilon_s}(x_n - x) & 0 \le x \le x_n \end{cases}
\end{equation}
Equating peak field at metallurgical junction $x=0$:
\begin{equation}
E(0) = -\frac{q N_A x_p}{\epsilon_s} = -\frac{q N_D x_n}{\epsilon_s} \implies \mathbf{N_A x_p = N_D x_n} \quad (\text{Charge Neutrality})
\end{equation}

\textbf{3. Integrating for Built-in Potential $V_{bi}$}:
Total potential drop across the depletion region equals the area under the triangular $E(x)$ field profile:
\begin{equation}
V_{bi} = -\int_{-x_p}^{x_n} E(x) dx = \frac{1}{2}|E_{\text{max}}| W = \frac{1}{2}\left(\frac{q N_D x_n}{\epsilon_s}\right)(x_p + x_n)
\end{equation}
Expressing in terms of carrier concentrations:
\begin{equation}
\mathbf{V_{bi} = \frac{k_BT}{q}\ln\left(\frac{N_A N_D}{n_i^2}\right)}
\end{equation}

\textbf{4. Total Depletion Width $W = x_p + x_n$}:
Using $x_p = W \frac{N_D}{N_A + N_D}$ and $x_n = W \frac{N_A}{N_A + N_D}$:
\begin{equation}
V_{bi} = \frac{q}{2\epsilon_s}\frac{N_A N_D}{N_A + N_D} W^2
\end{equation}
Under applied bias $V_a$ (forward $V_a > 0$, reverse $V_a = -V_R$):
\begin{equation}
\mathbf{W = \sqrt{\frac{2\epsilon_s}{q}\left(\frac{N_A + N_D}{N_A N_D}\right)(V_{bi} - V_a)} = \sqrt{\frac{2\epsilon_s}{q}\left(\frac{1}{N_A} + \frac{1}{N_D}\right)(V_{bi} - V_a)}}
\end{equation}
\end{theorysol}
\vspace{4mm}

\begin{theoryq}{6}
Describe the construction, working principle, energy band diagram, and $I$-$V$ characteristics of a Solar Cell. Define Open-Circuit Voltage ($V_{oc}$), Short-Circuit Current ($I_{sc}$), Fill Factor ($FF$), and Conversion Efficiency ($\eta$).
\end{theoryq}
\begin{theorysol}
\textbf{1. Construction}:
A solar cell is a large-area semiconductor p-n junction diode.
\begin{itemize}
    \item \textbf{Substrate}: Thick p-type Silicon wafer ($\sim 200\text{--}300\,\mu\text{m}$).
    \item \textbf{Emitter Layer}: Ultra-thin heavily doped n-type surface layer ($\sim 0.2\text{--}0.5\,\mu\text{m}$) to allow sunlight penetration.
    \item \textbf{Anti-Reflective Coating}: Thin film of $\text{Si}_3\text{N}_4$ ($n \approx 1.9$) to minimize surface reflection.
    \item \textbf{Contacts}: Metallic grid fingers on front surface (minimizing shading) and continuous solid metal contact on back.
\end{itemize}

\textbf{2. Working Principle}:
\begin{enumerate}
    \item \textbf{Photogeneration}: Sunlight with photon energy $h\nu \ge E_g$ penetrates the junction, exciting valence electrons into the conduction band and generating electron-hole pairs (EHPs).
    \item \textbf{Carrier Separation}: The strong built-in electric field in the space-charge region sweeps photo-generated electrons to the n-side and holes to the p-side before recombination can occur.
    \item \textbf{Power Delivery}: Carrier accumulation establishes an open-circuit photovoltage $V_{oc}$; connecting an external load resistor $R_L$ drives a continuous DC photocurrent.
\end{enumerate}

\textbf{3. Characteristic Parameters}:
\begin{enumerate}
    \item \textbf{Short-Circuit Current ($I_{sc}$)}: The maximum current when terminals are shorted ($V=0$):
    \begin{equation}
    I_{sc} \approx I_L = q A G (L_n + L_p + W) \propto P_{\text{incident}}
    \end{equation}
    \item \textbf{Open-Circuit Voltage ($V_{oc}$)}: The maximum voltage developed across open terminals ($I=0$):
    \begin{equation}
    \mathbf{V_{oc} = \frac{k_B T}{q}\ln\left(\frac{I_{sc}}{I_0} + 1\right)}
    \end{equation}
    \item \textbf{Fill Factor ($FF$)}: Ratio of maximum power rectangle to the product $V_{oc} I_{sc}$:
    \begin{equation}
    \mathbf{FF = \frac{P_{\text{max}}}{V_{oc} I_{sc}} = \frac{V_m I_m}{V_{oc} I_{sc}}}
    \end{equation}
    \item \textbf{Power Conversion Efficiency ($\eta$)}:
    \begin{equation}
    \mathbf{\eta = \frac{P_{\text{max}}}{P_{\text{in}}} = \frac{V_{oc} I_{sc} FF}{P_{\text{in}}} \times 100\%}
    \end{equation}
\end{enumerate}
\end{theorysol}
\vspace{4mm}

\begin{theoryq}{7}
Explain carrier transport in semiconductors under combined drift and diffusion. Derive the total electron and hole current density equations. State the continuity equation for excess minority carriers.
\end{theoryq}
\begin{theorysol}
\textbf{1. Drift Current}:
Drift current is driven by an applied electric field $\vec{E}$:
\begin{align}
\vec{J}_{n,\text{drift}} &= q n \mu_n \vec{E} \\
\vec{J}_{p,\text{drift}} &= q p \mu_p \vec{E}
\end{align}

\textbf{2. Diffusion Current}:
Diffusion current is driven by spatial carrier concentration gradients according to Fick's first law:
\begin{align}
\vec{J}_{n,\text{diff}} &= +q D_n \nabla n \\
\vec{J}_{p,\text{diff}} &= -q D_p \nabla p
\end{align}
where $D_n, D_p$ are electron and hole diffusion coefficients.

\textbf{3. Total Current Densities}:
Summing drift and diffusion components:
\begin{align}
\mathbf{\vec{J}_n = q n \mu_n \vec{E} + q D_n \nabla n} \\
\mathbf{\vec{J}_p = q p \mu_p \vec{E} - q D_p \nabla p}
\end{align}
Total conduction current density is $\vec{J}_{\text{total}} = \vec{J}_n + \vec{J}_p$.

\textbf{4. Continuity Equation for Excess Carriers}:
Applying conservation of charge to differential volume $A dx$:
\begin{align}
\mathbf{\frac{\partial (\Delta n)}{\partial t} = D_n \frac{\partial^2 (\Delta n)}{\partial x^2} + \mu_n E \frac{\partial (\Delta n)}{\partial x} + G_n - \frac{\Delta n}{\tau_n}} \\
\mathbf{\frac{\partial (\Delta p)}{\partial t} = D_p \frac{\partial^2 (\Delta p)}{\partial x^2} - \mu_p E \frac{\partial (\Delta p)}{\partial x} + G_p - \frac{\Delta p}{\tau_p}}
\end{align}
where $G$ is external generation rate and $\tau$ is minority carrier recombination lifetime.
\end{theorysol}
\vspace{4mm}

\begin{theoryq}{8}
Discuss the temperature dependence of carrier concentration in an extrinsic n-type semiconductor. Explain the three distinct operational regimes: (a) Freeze-out regime, (b) Extrinsic saturation regime, and (c) Intrinsic transition regime.
\end{theoryq}
\begin{theorysol}
\textbf{1. Temperature Variation of Electron Concentration $n(T)$}:
A graph of $\ln(n)$ versus $1/T$ exhibits three distinct physical regimes:

\begin{enumerate}
    \item \textbf{Low-Temperature Freeze-Out Regime ($T < 100\text{ K}$)}:
    Thermal energy is lower than the donor ionization energy ($k_BT \ll E_c - E_d \approx 0.045\text{ eV}$).
    Electrons lack sufficient thermal energy to jump into the conduction band and freeze back into localized donor states.
    \begin{equation}
    n(T) \approx \sqrt{\frac{N_c N_D}{2}} \exp\left(-\frac{E_c - E_d}{2 k_B T}\right) \propto e^{-\Delta E_d / 2k_BT}
    \end{equation}
    Carrier concentration increases exponentially with rising temperature.

    \item \textbf{Moderate-Temperature Extrinsic (Saturation/Exhaustion) Regime ($100\text{ K} \le T \le 450\text{ K}$)}:
    Thermal energy ($k_BT \ge 26\text{ meV}$) is sufficient to ionize $100\%$ of all donor impurity atoms ($N_D^+ \approx N_D$), while intrinsic band-to-band thermal generation across the large bandgap ($E_g = 1.12\text{ eV}$) is still negligible ($n_i \ll N_D$).
    \begin{equation}
    \mathbf{n \approx N_D = \text{constant}}
    \end{equation}
    This is the stable operating regime for all standard silicon electronic devices.

    \item \textbf{High-Temperature Intrinsic Regime ($T > 500\text{ K}$)}:
    Thermal generation of electron-hole pairs across the main forbidden gap increases exponentially ($n_i(T) \propto T^{3/2}e^{-E_g/2k_BT}$) and quickly overwhelms the fixed dopant concentration ($n_i(T) \gg N_D$).
    \begin{equation}
    \mathbf{n \approx n_i(T) \propto e^{-E_g / 2k_BT}}
    \end{equation}
    The semiconductor loses its extrinsic properties and behaves as an intrinsic material, setting the upper thermal operating limit of silicon chips.
\end{enumerate}
\end{theorysol}
\vspace{4mm}

\begin{theoryq}{9}
Explain the Direct vs Indirect bandgap semiconductor concepts using $E$-$k$ energy-momentum diagrams. Discuss the optoelectronic selection rules and why indirect semiconductors cannot be used for efficient lasers or LEDs.
\end{theoryq}
\begin{theorysol}
\textbf{1. $E$-$k$ Diagram \& Bandgap Classification}:
In solid-state physics, the energy states of electrons are plotted as parabolic bands $E(k)$ versus crystal momentum wavevector $k = p/\hbar$.
\begin{itemize}
    \item \textbf{Direct Bandgap Semiconductor (e.g., GaAs, InP, GaN)}:
    The conduction band minimum ($E_c$) and valence band maximum ($E_v$) align at the exact same wavevector coordinate ($k = 0$, $\Gamma$-point).
    \item \textbf{Indirect Bandgap Semiconductor (e.g., Si, Ge, GaP)}:
    The valence band maximum is at $k = 0$, but the conduction band minimum is located at an off-axis wavevector $k = k_0 \ne 0$ along the $(100)$ crystal direction.
\end{itemize}

\textbf{2. Conservation Laws for Optical Transitions}:
A photon carrying energy $h\nu \approx E_g \sim 1\text{ eV}$ has wavevector $k_{\text{photon}} = \frac{2\pi}{\lambda} \approx 10^7\text{ m}^{-1}$.
By contrast, the Brillouin zone boundary has $k_{\text{crystal}} = \pi/a \approx 10^{10}\text{ m}^{-1}$, which is $1000\times$ larger ($k_{\text{photon}} \approx 0$).

\begin{enumerate}
    \item \textbf{Direct Bandgap Recombination}:
    Because $\Delta k = 0$, an electron at the conduction band bottom recombines vertically with a hole at the valence band top. Momentum is conserved directly by the photon alone:
    \begin{equation}
    \mathbf{e^- + h^+ \to \text{Photon } (h\nu)} \quad (\text{Fast, 1st-order process: } \tau_{\text{rad}} \sim 1\text{ ns})
    \end{equation}
    Quantum radiative efficiency approaches $100\%$.

    \item \textbf{Indirect Bandgap Recombination}:
    Because $\Delta k = k_0 \ne 0$, an electron cannot drop vertically to the valence top without transferring substantial momentum to the crystal lattice.
    Recombination requires simultaneous emission or absorption of a lattice vibration quantum (\textbf{Phonon} $\hbar\Omega$):
    \begin{equation}
    \mathbf{e^- + h^+ \to \text{Photon } (h\nu) \pm \text{Phonon } (\hbar\Omega)} \quad (\text{Slow, 2nd-order 3-body process: } \tau_{\text{rad}} \sim 1\text{ ms})
    \end{equation}
    Because this process is $10^6\times$ slower, excited carriers undergo non-radiative defect recombination, releasing energy as useless heat rather than coherent light.
\end{enumerate}
\end{theorysol}
\vspace{4mm}

\begin{theoryq}{10}
A Hall effect measurement is performed on an unknown semiconductor slab of length $L = 1.0\text{ cm}$, width $w = 2.0\text{ mm}$, and thickness $d = 0.4\text{ mm}$. A current of $I = 15\text{ mA}$ is passed along the length in a magnetic field $B = 0.8\text{ T}$ applied along the thickness. A Hall voltage $V_H = -6.4\text{ mV}$ is measured across the width, and the voltage drop along the length is $V_L = 1.2\text{ V}$. Determine: (a) carrier type, (b) Hall coefficient $R_H$, (c) majority carrier concentration $n$, (d) electrical resistivity $\rho$ and conductivity $\sigma$, and (e) Hall mobility $\mu_H$.
\end{theoryq}
\begin{theorysol}
\textbf{1. Carrier Type}:
Because the measured Hall voltage is negative ($V_H = -6.4\text{ mV} < 0$), majority carriers are electrons. The sample is confirmed to be an \textbf{n-type semiconductor}.

\textbf{2. Hall Coefficient $R_H$}:
Thickness $d = 0.4\text{ mm} = 4.0 \times 10^{-4}\text{ m}$.
\begin{equation}
\mathbf{R_H = \frac{V_H d}{I B} = \frac{(-6.4 \times 10^{-3}\text{ V})(4.0 \times 10^{-4}\text{ m})}{(15 \times 10^{-3}\text{ A})(0.8\text{ T})} = \frac{-2.56 \times 10^{-6}}{0.012} = -2.133 \times 10^{-4}\text{ m}^3/\text{C}}
\end{equation}

\textbf{3. Carrier Concentration $n$}:
\begin{equation}
\mathbf{n = \frac{1}{|R_H| e} = \frac{1}{(2.133 \times 10^{-4}\text{ m}^3/\text{C})(1.602 \times 10^{-19}\text{ C})} \approx 2.926 \times 10^{22}\text{ m}^{-3} = 2.93 \times 10^{16}\text{ cm}^{-3}}
\end{equation}

\textbf{4. Electrical Resistivity $\rho$ \& Conductivity $\sigma$}:
Cross-sectional area $A = w \times d = (2 \times 10^{-3}\text{ m})(4 \times 10^{-4}\text{ m}) = 8 \times 10^{-7}\text{ m}^2$.
Sample resistance $R = V_L / I = 1.2\text{ V} / (15 \times 10^{-3}\text{ A}) = 80\,\Omega$.
\begin{align*}
\mathbf{\rho} &= \frac{R A}{L} = \frac{(80\,\Omega)(8 \times 10^{-7}\text{ m}^2)}{1.0 \times 10^{-2}\text{ m}} = \mathbf{6.4 \times 10^{-3}\,\Omega\cdot\text{m}} \\
\mathbf{\sigma} &= \frac{1}{\rho} = \frac{1}{6.4 \times 10^{-3}} = \mathbf{156.25\text{ S/m}}
\end{align*}

\textbf{5. Hall Mobility $\mu_H$}:
\begin{align*}
\mathbf{\mu_H} &= |R_H| \sigma = (2.133 \times 10^{-4}\text{ m}^3/\text{C})(156.25\text{ S/m}) \\
&= \mathbf{0.0333\text{ m}^2/\text{V}\cdot\text{s} = 333.3\text{ cm}^2/\text{V}\cdot\text{s}}
\end{align*}
\end{theorysol}
\vspace{4mm}

\begin{theoryq}{11}
Derive the temperature dependence of the Fermi level $E_F(T)$ in an n-type extrinsic semiconductor. Explain the physical behavior in the cryogenic freeze-out, extrinsic saturation, and high-temperature intrinsic regimes.
\end{theoryq}
\begin{theorysol}
\textbf{1. Charge Neutrality Condition}:
In an n-type semiconductor with donor concentration $N_d$ (energy level $E_d$) and no acceptors:
\begin{equation}
n = N_d^+ + p
\end{equation}
For non-degenerate electrons: $n = N_c e^{-(E_c - E_F)/k_B T}$, $p = N_v e^{-(E_F - E_v)/k_B T}$, and ionized donors $N_d^+ = \frac{N_d}{1 + 2 e^{(E_F - E_d)/k_B T}}$.

\textbf{2. Three Distinct Temperature Regimes}:
\begin{enumerate}
\item \textbf{Freeze-Out / Ionization Regime ($T \to 0\text{ K}$):}
Thermal energy is insufficient to ionize donors. Here $p \approx 0$ and $n = N_d^+$.
\begin{equation}
N_c e^{-(E_c - E_F)/k_B T} \approx N_d e^{-(E_F - E_d)/k_B T} \implies \mathbf{E_F(0) = \frac{E_c + E_d}{2}}
\end{equation}
As $T$ rises slightly, $E_F(T) = \frac{E_c + E_d}{2} + \frac{k_B T}{2}\ln\left(\frac{N_d}{2 N_c}\right)$.

\item \textbf{Extrinsic / Saturation Regime ($100\text{ K} \le T \le 450\text{ K}$):}
All donor atoms are completely ionized ($N_d^+ \approx N_d$) and minority carrier generation is negligible ($p \ll n$):
\begin{equation}
n = N_d \implies N_c e^{-(E_c - E_F)/k_B T} = N_d
\end{equation}
\begin{equation}
\mathbf{E_F(T) = E_c - k_B T \ln\left(\frac{N_c}{N_d}\right)}
\end{equation}
In this range, $n = N_d = \text{constant}$, and the Fermi level decreases linearly with temperature toward the middle of the bandgap.

\item \textbf{Intrinsic Crossover Regime ($T > 500\text{ K}$):}
Band-to-band thermal generation dominates over donor doping ($n_i(T) \gg N_d$):
\begin{equation}
n \approx p \approx n_i \implies \mathbf{E_F \to E_i = \frac{E_c + E_v}{2} + \frac{3}{4}k_B T \ln\left(\frac{m_h^*}{m_e^*}\right)}
\end{equation}
The semiconductor loses its extrinsic characteristics and behaves as an intrinsic semiconductor.
\end{enumerate}
\end{theorysol}
\vspace{4mm}

\begin{theoryq}{12}
Compare direct and indirect carrier recombination mechanisms in semiconductors. Explain Shockley-Read-Hall (SRH) recombination through deep trap levels, Auger recombination, and radiative recombination.
\end{theoryq}
\begin{theorysol}
\textbf{1. Direct vs Indirect Bandgap Recombination}:
\begin{itemize}
\item \textbf{Direct Recombination (e.g., GaAs, InP):} Electron at conduction band minimum drops directly into valence band maximum ($\Delta k = 0$). Momentum is conserved without phonon participation; spontaneous recombination lifetime is short ($\tau_r \sim 1\text{--}10\text{ ns}$), emitting photons efficiently.
\item \textbf{Indirect Recombination (e.g., Si, Ge):} Conduction band minimum and valence band maximum are offset in $k$-space ($\Delta k \ne 0$). Direct photon emission is forbidden by momentum conservation; transition requires simultaneous emission/absorption of a lattice phonon, leading to very long radiative lifetimes ($\tau_r \sim 1\text{ ms}$).
\end{itemize}

\textbf{2. Shockley-Read-Hall (SRH) Trap-Assisted Recombination}:
In indirect semiconductors like Silicon, non-radiative recombination occurs primarily via deep-level defect/impurity states located near the center of the bandgap ($E_t \approx E_i$):
\begin{enumerate}
\item Trap captures an electron from the conduction band.
\item Trap captures a hole from the valence band (or emits trapped electron to valence band).
\item Electron and hole annihilate, dissipating energy into the lattice as heat (phonons).
\end{enumerate}
The SRH recombination rate is:
\begin{equation}
U_{SRH} = \frac{np - n_i^2}{\tau_{p0}(n + n_1) + \tau_{n0}(p + p_1)}
\end{equation}

\textbf{3. Auger Recombination}:
A three-particle non-radiative process where the recombination energy of an electron-hole pair is transferred to a third carrier (electron or hole), which relaxes thermally.
The Auger recombination rate scales as $U_{Auger} = C_n n^2 p + C_p n p^2$, dominating at very high injection carrier densities ($n > 10^{18}\text{ cm}^{-3}$).
\end{theorysol}
\vspace{4mm}

\begin{theoryq}{13}
Explain carrier velocity saturation at high electric fields in Silicon and describe the Gunn Effect / Negative Differential Resistance (NDR) in Gallium Arsenide (GaAs) using the Ridley-Watkins-Hilsum two-valley model.
\end{theoryq}
\begin{theorysol}
\textbf{1. High-Field Velocity Saturation}:
At low electric fields, drift velocity increases linearly with field according to Ohm's law: $v_d = \mu E$.
At high electric fields ($E > 10^4\text{ V/cm}$), carriers gain kinetic energy faster than they can transfer it to acoustic phonons. Optical phonon emission becomes intense, scattering carriers randomly and saturating the drift velocity to a constant value:
\begin{equation}
v_{sat} \approx 1.0 \times 10^7\text{ cm/s}\quad (\text{in Silicon at }300\text{ K})
\end{equation}
The empirical drift velocity relation is:
\begin{equation}
v_d(E) = \frac{\mu_0 E}{\left[1 + \left(\frac{\mu_0 E}{v_{sat}}\right)^\beta\right]^{1/\beta}}
\end{equation}

\textbf{2. Gunn Effect and Negative Differential Resistance in GaAs}:
In GaAs, the conduction band has a two-valley structure:
\begin{itemize}
\item \textbf{Lower Valley ($\Gamma$-valley):} Located at $k=0$; low effective mass ($m_1^* = 0.067 m_0$), high electron mobility ($\mu_1 \approx 8500\text{ cm}^2/\text{V}\cdot\text{s}$).
\item \textbf{Upper Satellite Valleys ($L$-valleys):} Separated by energy offset $\Delta E \approx 0.31\text{ eV}$; high effective mass ($m_2^* = 0.22 m_0$), low mobility ($\mu_2 \approx 100\text{ cm}^2/\text{V}\cdot\text{s}$).
\end{itemize}

\textbf{3. Mechanism of Negative Differential Mobility}:
When applied electric field exceeds the threshold field $E_{th} \approx 3.2\text{ kV/cm}$:
\begin{enumerate}
\item Electrons gain sufficient kinetic energy to transfer from the high-mobility $\Gamma$-valley into the low-mobility $L$-valleys (intervalley scattering).
\item Average electron drift velocity $v_d = \frac{n_1 \mu_1 + n_2 \mu_2}{n_1 + n_2} E$ decreases as field increases above $E_{th}$.
\item This yields a negative differential mobility regime:
\begin{equation}
\mu_d = \frac{dv_d}{dE} < 0
\end{equation}
\end{enumerate}
This instability produces moving high-field charge domains (Gunn domains), generating microwave oscillations ($1\text{--}100\text{ GHz}$) in Gunn diodes.
\end{theorysol}
\vspace{4mm}

\begin{theoryq}{14}
Derive the non-ideal diode I-V equation including recombination in the depletion region. Explain the physical origin of the Diode Ideality Factor $\eta$ ($\eta=1$ vs $\eta=2$) and series resistance effects.
\end{theoryq}
\begin{theorysol}
\textbf{1. Total Diode Current Components}:
The total forward current in a real p-n junction consists of:
\begin{equation}
I = I_{diff} + I_{rec}
\end{equation}
where $I_{diff}$ is the ideal diffusion current in quasi-neutral regions and $I_{rec}$ is the carrier recombination current within the depletion space-charge layer.

\textbf{2. Depletion Region Recombination Current}:
Inside the depletion region of width $W$, the SRH recombination rate reaches maximum when $n \approx p \approx n_i e^{qV / (2k_B T)}$.
Integrating $U_{SRH}$ across the depletion width $W$:
\begin{equation}
I_{rec} = q A \int_0^W U_{SRH}\, dx \approx \frac{q A W n_i}{2\tau_0} e^{q V / (2k_B T)} = I_{r0} e^{q V / (2k_B T)}
\end{equation}

\textbf{3. Generalized Diode Equation and Ideality Factor $\eta$}:
Combining both components yields the empirical Shockley diode equation:
\begin{equation}
\mathbf{I = I_0 \left[\exp\left(\frac{q V}{\eta k_B T}\right) - 1\right]}
\end{equation}
\begin{itemize}
\item \textbf{At Low Forward Bias ($V < 0.3\text{ V}$):} Recombination in the depletion layer dominates $\implies \mathbf{\eta \approx 2}$.
\item \textbf{At Intermediate Forward Bias ($0.3\text{ V} < V < 0.7\text{ V}$):} Diffusion of injected minority carriers in neutral regions dominates $\implies \mathbf{\eta \approx 1}$.
\item \textbf{At High Forward Bias ($V > 0.7\text{ V}$):} Bulk ohmic series resistance $R_s$ causes the external voltage to drop as $V_{internal} = V_{ext} - I R_s$, bending the $\ln(I)$ vs $V$ curve downward.
\end{itemize}

\end{theorysol}
\vspace{4mm}

\begin{theoryq}{15}
Explain the photovoltaic effect in a solar cell. Derive the expressions for Short-Circuit Current $I_{sc}$, Open-Circuit Voltage $V_{oc}$, Maximum Power Output $P_{max}$, Fill Factor ($FF$), and Power Conversion Efficiency $\eta$.
\end{theoryq}
\begin{theorysol}
\textbf{1. Solar Cell Working Principle}:
When incident photons with $h\nu \ge E_g$ strike the p-n junction, electron-hole pairs are photogenerated. The built-in electric field $\mathbf{E}_{bi}$ in the depletion region separates these carriers, sweeping electrons to the n-side and holes to the p-side, establishing a photovoltage.

\textbf{2. Equivalent Circuit and Current Equation}:
The total current delivered to a load under illumination is:
\begin{equation}
I(V) = I_L - I_0 \left(e^{\frac{qV}{k_B T}} - 1\right)
\end{equation}
where $I_L$ is the photogenerated light current and $I_0$ is the reverse saturation current.

\textbf{3. Key Solar Cell Parameters}:
\begin{itemize}
\item \textbf{Short-Circuit Current ($I_{sc}$):} Setting $V=0$:
\begin{equation}
\mathbf{I_{sc} = I_L}
\end{equation}
\item \textbf{Open-Circuit Voltage ($V_{oc}$):} Setting $I = 0$:
\begin{equation}
0 = I_L - I_0 \left(e^{\frac{qV_{oc}}{k_B T}} - 1\right) \implies \mathbf{V_{oc} = \frac{k_B T}{q} \ln\left(\frac{I_L}{I_0} + 1\right) \approx \frac{k_B T}{q}\ln\left(\frac{I_{sc}}{I_0}\right)}
\end{equation}
\item \textbf{Maximum Power Point ($P_{max}$):} Power output $P(V) = V I(V)$ is maximized when $\frac{dP}{dV} = 0$:
\begin{equation}
P_{max} = V_{mp} I_{mp}
\end{equation}
\item \textbf{Fill Factor ($FF$):} Represents the squareness of the I-V curve:
\begin{equation}
\mathbf{FF = \frac{P_{max}}{V_{oc} I_{sc}} = \frac{V_{mp} I_{mp}}{V_{oc} I_{sc}}}\quad (\text{Typically }0.75\text{--}0.85)
\end{equation}
\item \textbf{Power Conversion Efficiency ($\eta_{cell}$):}
\begin{equation}
\mathbf{\eta = \frac{P_{max}}{P_{in}} = \frac{V_{oc} I_{sc} FF}{P_{in} \cdot A_{cell}} \times 100\%}
\end{equation}
where $P_{in}$ is the incident optical irradiance ($1000\text{ W/m}^2$ under standard AM1.5G test conditions).
\end{itemize}

\end{theorysol}
\vspace{4mm}

\begin{theoryq}{16}
Explain the physics of Light Emitting Diodes (LEDs). Derive the internal quantum efficiency $\eta_{int}$, external quantum efficiency $\eta_{ext}$, the critical extraction angle $\theta_c$, and methods used to enhance light extraction.
\end{theoryq}
\begin{theorysol}
\textbf{1. Principle of Operation}:
Under forward bias, minority carriers are injected across the p-n junction. Spontaneous radiative recombination of injected electrons and holes releases energy as photons of wavelength:
\begin{equation}
\lambda \approx \frac{h c}{E_g} = \frac{1.24}{E_g\text{ [eV]}}\; \mu\text{m}
\end{equation}

\textbf{2. Internal Quantum Efficiency ($\eta_{int}$)}:
Defined as the ratio of radiative recombination rate ($R_r$) to total recombination rate ($R_r + R_{nr}$):
\begin{equation}
\mathbf{\eta_{int} = \frac{R_r}{R_r + R_{nr}} = \frac{\tau_{nr}}{\tau_r + \tau_{nr}}}
\end{equation}
In direct bandgap semiconductors (InGaN, AlGaInP), $\tau_r \ll \tau_{nr}$, making $\eta_{int} > 80\%$.

\textbf{3. External Quantum Efficiency and Light Extraction Cone}:
Semiconductors have high refractive indices (e.g., $n_{GaAs} \approx 3.6$).
By Snell's law, total internal reflection traps all light striking the surface at angles greater than the critical angle $\theta_c$:
\begin{equation}
\theta_c = \sin^{-1}\left(\frac{n_{air}}{n_{semi}}\right) = \sin^{-1}\left(\frac{1}{3.6}\right) \approx 16.1^\circ
\end{equation}
The fraction of light escaping within the escape cone from an isotropic internal point source is:
\begin{equation}
\eta_{escape} \approx \frac{\Omega_{cone}}{4\pi} = \frac{2\pi(1 - \cos\theta_c)}{4\pi} \approx \frac{1 - \sqrt{1 - (1/n_s)^2}}{2} \approx \frac{1}{4 n_s^2}
\end{equation}
For GaAs ($n=3.6$), $\eta_{escape} \approx \frac{1}{4(3.6)^2} \approx 1.93\%$.

\textbf{4. Techniques to Enhance Extraction Efficiency}:
\begin{itemize}
\item Transparent hemispherical epoxy encapsulation ($n_{epoxy} \approx 1.5$) widens escape cone.
\item Surface texturing / roughening creates multiple scattering angles for trapped photons.
\item Resonant cavity and photonic crystal patterns redirect trapped waveguide modes outward.
\end{itemize}

\end{theorysol}
\vspace{4mm}

\begin{theoryq}{17}
Describe the working principle of Avalanche Photodiodes (APDs). Derive the expression for optical responsivity $\mathcal{R}$, quantum efficiency $\eta_q$, and explain impact ionization and avalanche multiplication factor $M$.
\end{theoryq}
\begin{theorysol}
\textbf{1. Photodetector Operation and Optical Absorption}:
Incident photons with energy $h\nu \ge E_g$ generate electron-hole pairs in the high-field reverse-biased depletion region.
The primary photocurrent is:
\begin{equation}
I_{ph0} = q \eta_q \left(\frac{P_{opt}}{h\nu}\right)
\end{equation}
where $\eta_q$ is the quantum efficiency ($\eta_q = (1 - R)(1 - e^{-\alpha W})$) and $P_{opt}$ is incident optical power.

\textbf{2. Optical Responsivity ($\mathcal{R}$)}:
The responsivity in Amperes per Watt is:
\begin{equation}
\mathbf{\mathcal{R} = \frac{I_{ph}}{P_{opt}} = \frac{q \eta_q}{h\nu} M = \frac{q \eta_q \lambda}{h c} M \approx \frac{\eta_q \lambda\text{ [}\mu\text{m]}}{1.24} M\quad [\text{A/W}]}
\end{equation}
where $M$ is the avalanche multiplication factor ($M=1$ for PIN photodiodes).

\textbf{3. Impact Ionization and Avalanche Multiplication}:
In an APD, a high reverse bias ($>100\text{ V}$) produces an internal electric field exceeding $10^5\text{ V/cm}$.
\begin{enumerate}
\item Injected primary photoelectrons are accelerated to very high kinetic energies.
\item High-energy electrons collide with valence electrons in the lattice, knocking them into the conduction band to create secondary electron-hole pairs (\textit{impact ionization}).
\item The secondary carriers are also accelerated, triggering a cascade chain reaction.
\end{enumerate}

\textbf{4. Multiplication Factor $M$ (Miller's Empirical Formula)}:
\begin{equation}
\mathbf{M = \frac{1}{1 - \left(\frac{V_R}{V_B}\right)^m}}
\end{equation}
where $V_B$ is the avalanche breakdown voltage and $m \approx 3\text{--}6$.
APDs provide internal current gains of $M \sim 50\text{--}200$, significantly enhancing receiver sensitivity in optical communication links.
\end{theorysol}
\vspace{4mm}

\begin{theoryq}{18}
Explain the structure and operational principles of the MOSFET (Metal-Oxide-Semiconductor Field-Effect Transistor). Draw the MOS energy band diagrams under Accumulation, Depletion, and Strong Inversion, and derive the Threshold Voltage $V_{th}$.
\end{theoryq}
\begin{theorysol}
\textbf{1. MOS Capacitor Structure}:
Consists of a metal/polysilicon gate, a thin dielectric oxide layer ($\text{SiO}_2$, thickness $t_{ox}$, capacitance per unit area $C_{ox} = \epsilon_{ox}/t_{ox}$), and a p-type silicon substrate.

\textbf{2. Three Operational Regimes and Energy Band Diagrams}:
\begin{itemize}
\item \textbf{Accumulation ($V_G < V_{FB}$):} Negative gate voltage attracts majority holes to the $\text{Si-SiO}_2$ interface. Valence band bends upward toward Fermi level.
\item \textbf{Depletion ($V_{FB} < V_G < V_{th}$):} Positive gate bias repels holes away from the surface, creating an uncompensated negative acceptor space charge region ($N_a^-$) of width $W_{dep} = \sqrt{\frac{2\epsilon_s \psi_s}{q N_a}}$. Bands bend downward; surface potential $0 < \psi_s < 2\phi_B$.
\item \textbf{Strong Inversion ($V_G \ge V_{th}$):} Bands bend downward such that the surface potential reaches:
\begin{equation}
\psi_s = 2\phi_B = 2\left(\frac{k_B T}{q}\ln\frac{N_a}{n_i}\right)
\end{equation}
The conduction band approaches the Fermi level, attracting a high concentration of free electrons to form an n-type conducting inversion channel between Source and Drain.
\end{itemize}

\textbf{3. Threshold Voltage Derivation}:
The gate voltage required to achieve strong inversion ($V_{th}$) equals flatband voltage plus oxide voltage drop plus surface potential:
\begin{equation}
\mathbf{V_{th} = V_{FB} + 2\phi_B + \frac{\sqrt{2\epsilon_s q N_a (2\phi_B)}}{C_{ox}}}
\end{equation}
where $V_{FB} = \Phi_{ms} - \frac{Q_{ox}}{C_{ox}}$ accounts for work function difference and oxide charges.
\end{theorysol}
\vspace{4mm}

\begin{theoryq}{19}
Compare Zener Breakdown and Avalanche Breakdown in reverse-biased p-n junctions. Contrast their physical mechanisms, breakdown voltage thresholds, electric field values, and temperature coefficients.
\end{theoryq}
\begin{theorysol}
\textbf{1. Comprehensive Comparison Table}:
\begin{center}
\begin{tabularx}{\textwidth}{|l|X|X|}
\hline
\textbf{Parameter} & \textbf{Zener Breakdown} & \textbf{Avalanche Breakdown} \\
\hline
\textbf{Doping Concentration} & Very heavily doped ($N_a, N_d > 10^{18}\text{ cm}^{-3}$) & Lightly to moderately doped ($N_a, N_d < 10^{17}\text{ cm}^{-3}$) \\
\textbf{Depletion Width} & Extremely thin ($W < 10\text{ nm} = 100\text{ \AA}$) & Wide ($W > 1\;\mu\text{m}$) \\
\textbf{Electric Field} & Intense ($E > 10^6\text{ V/cm} = 10^8\text{ V/m}$) & Moderate ($E \sim 10^5\text{ V/cm}$) \\
\textbf{Physical Mechanism} & Quantum mechanical tunneling of valence electrons across narrow gap & Impact ionization / carrier multiplication by energetic carriers \\
\textbf{Breakdown Voltage} & Low ($V_Z < 5\text{--}6\text{ V}$) & High ($V_B > 6\text{--}8\text{ V}$) \\
\textbf{Temperature Coefficient} & \textbf{Negative} ($\frac{dV_Z}{dT} < 0$): $E_g$ decreases with $T$, making tunneling easier & \textbf{Positive} ($\frac{dV_B}{dT} > 0$): Lattice vibrations increase phonon scattering, requiring higher field \\
\textbf{I-V Knee Characteristic} & Very sharp and stable & Slightly softer knee \\
\hline
\end{tabularx}
\end{center}

\textbf{2. Practical Implication}:
Zener diodes engineered with breakdown voltage near $\sim 5.6\text{ V}$ combine Zener and Avalanche mechanisms with opposite temperature coefficients, achieving a near-zero overall temperature drift ideal for precision voltage references.
\end{theorysol}
\vspace{4mm}

\begin{theoryq}{20}
A Hall effect measurement is performed on a rectangular slice of Silicon with dimensions $L = 10\text{ mm}$, width $w = 2.0\text{ mm}$, and thickness $t = 0.5\text{ mm}$. A longitudinal current $I_x = 5.0\text{ mA}$ is applied. In a magnetic field $B_z = 0.60\text{ T}$, the measured Hall voltage is $V_H = -12.5\text{ mV}$. The longitudinal voltage drop along length $L$ is $V_x = 2.0\text{ V}$. (a) Determine the majority carrier type and concentration $n$. (b) Calculate the Hall coefficient $R_H$. (c) Find the electrical conductivity $\sigma$ and electron drift mobility $\mu_n$.
\end{theoryq}
\begin{theorysol}
\textbf{1. Determination of Majority Carrier Type}:
The measured Hall voltage is negative ($V_H = -12.5\text{ mV}$), indicating that accumulated transverse carriers are electrons.
Therefore, the sample is an \textbf{n-type semiconductor}.

\textbf{2. Hall Coefficient ($R_H$) and Carrier Density ($n$)}:
From the Hall voltage formula:
\begin{equation}
V_H = \frac{R_H I_x B_z}{t} \implies \mathbf{R_H = \frac{V_H t}{I_x B_z}}
\end{equation}
\begin{align*}
R_H &= \frac{(-12.5 \times 10^{-3}\text{ V})(0.5 \times 10^{-3}\text{ m})}{(5.0 \times 10^{-3}\text{ A})(0.60\text{ T})} = \frac{-6.25 \times 10^{-6}}{3.0 \times 10^{-3}} = \mathbf{-2.083 \times 10^{-3}\text{ m}^3/\text{C}}
\end{align*}
Carrier concentration:
\begin{equation}
\mathbf{n} = \frac{1}{q |R_H|} = \frac{1}{(1.602 \times 10^{-19}\text{ C})(2.083 \times 10^{-3}\text{ m}^3/\text{C})} = \mathbf{3.00 \times 10^{21}\text{ m}^{-3}} = \mathbf{3.00 \times 10^{15}\text{ cm}^{-3}}
\end{equation}

\textbf{3. Electrical Conductivity ($\sigma$)}:
The sample cross-sectional area is $A = w \times t = (2.0 \times 10^{-3})(0.5 \times 10^{-3}) = 1.0 \times 10^{-6}\text{ m}^2$.
Sample resistance:
\begin{equation}
R = \frac{V_x}{I_x} = \frac{2.0\text{ V}}{5.0 \times 10^{-3}\text{ A}} = 400\text{ }\Omega
\end{equation}
Conductivity:
\begin{equation}
\mathbf{\sigma} = \frac{L}{R A} = \frac{10 \times 10^{-3}\text{ m}}{(400\text{ }\Omega)(1.0 \times 10^{-6}\text{ m}^2)} = \mathbf{25.0\text{ }\Omega^{-1}\text{m}^{-1}\text{ (or S/m)}}
\end{equation}

\textbf{4. Electron Drift Mobility ($\mu_n$)}:
\begin{equation}
\mathbf{\mu_n} = \sigma |R_H| = (25.0\text{ S/m})(2.083 \times 10^{-3}\text{ m}^3/\text{C}) = \mathbf{0.0521\text{ m}^2/(\text{V}\cdot\text{s})} = \mathbf{521\text{ cm}^2/(\text{V}\cdot\text{s})}
\end{equation}
\end{theorysol}
\vspace{4mm}
